\chapter{Stage Appendix: Part VII -- One-Port Same-Charge Sieve, Actual-Branch Tests, and Rigid-Mouth Orbit Lock}
\label{app:stage-part-7}

\section*{Stage range: 202--225}
\addcontentsline{toc}{section}{Stage range: 202--225}

This appendix expands the Part VII theorem-block chapter into a stage-by-stage verification ledger.  The main-body chapter states the same-charge and rigid-mouth conclusions in citation-ready form; the present appendix keeps the derivational path auditable, preserving the exact stage order from Stage~202 through Stage~225.  The stages are intentionally grouped into three verification questions: first, whether the one-port mixed sector creates a new long-range same-charge kernel; second, whether a short-range or resonant corridor survives the transported prefactor ceilings; and third, whether the actual rigid-mouth branch lands on the selected support and orbit-lock packet.

The appendix should be read with the same status firewall as the main body.  The static kernels, projectors, packet maps, and selected-support compilers are exact inside their declared closures.  Numerical pole censuses and known-data insertions are numerical placements into exact formulas.  The realized nonlinear moving-throat branch remains an open input wherever the text asks whether the actual PDE branch lands on a selected point, window, or orbit-lock packet.

Part VI closed the local realization compiler: once a branch is expressed in the free-quintuple graph coordinates, the remaining local mixed-ray search is finite through support-cardinality five.  Part VII is the first use of that completed back end on the same-charge, same-sign, one-port mixed-sector problem.  Its job is not to add another algebraic search layer.  Its job is to insert actual one-port wall/BdG/Maxwell/mixed branch data into the compiler, decide what the same-charge mixed sector can and cannot do, and then translate the surviving branch tests into the rigid-mouth physical normal form.

The result anchor served by this chapter is \resultanchor{MTDC-T11}.  This anchor should be cited when a downstream document needs the derivation of the one-port static same-charge kernel, the dynamic phase-lag and resonance/linewidth tests, the transported actual-branch prefactor ceiling, the strict first-order even-gate package, the pure-transfer/co-loading sieve, the selected-branch dynamic-window compiler, the rigid-mouth packet projectors, the physical logarithmic chart \((U,V)\), the Cartesian orbit-lock theorem, or the selected twin-support placement compiler.

The conceptual outcome is simple but important.  The mixed sector remains real and necessary, but it does not supply a new long-range same-charge attraction.  The static one-port bundle produces only short-range product kernels.  The linear dynamic bundle produces no new spatial kernel class and its passive/outgoing correction is phase lag at first order.  What survives is a sharply constrained short-range/resonant corridor and, on the actual rigid-mouth branch, a two-coordinate orbit-lock packet.  The rigid-mouth same-charge problem is diagonal in the physical logarithmic chart
\begin{equation}
\label{eq:app-part07-intro-uv-chart}
U:=\ln\frac{\mathcal T^2}{\mathcal T_{\rm ref}^2},
\qquad
V:=\ln\frac{\epsilon_\eta}{\epsilon_{\eta,\rm ref}},
\end{equation}
and the strict orbit-lock point is
\begin{equation}
\label{eq:app-part07-intro-orbit-lock}
U=0,
\qquad
V=0.
\end{equation}
The selected passive/outgoing support side is also no longer a three-regime ambiguity: the minimal isotropic quadrupole precursor fixes
\begin{equation}
\label{eq:app-part07-intro-selected-twin-slice}
\rho_\alpha=\frac43,
\qquad
\zeta_{\rm req}=\frac13,
\qquad
\Pi_{\rm tr}=\frac43 C_{\rm mix},
\end{equation}
placing the selected branch on the symmetric lowest-twin support slice.

\begin{longtable}{@{}L{0.075\textwidth}L{0.34\textwidth}L{0.20\textwidth}L{0.33\textwidth}@{}}
\caption{Part VII source and status map.}\label{tab:app-part07-source-status-map}\\
\toprule
Stage & Working title & Primary status & Main output \\
\midrule
\endfirsthead
\toprule
Stage & Working title & Primary status & Main output \\
\midrule
\endhead
202 & One-port mixed-bundle static kernel & \StatusExactClosure{} & Exact static \(3\times3\) bundle, determinant \(\det\mathcal K_{\rm red}=\Delta D_0\), susceptibility kernel, outgoing-prefactor bridge, and square-law suppression verdict. \\
\addlinespace
203 & Dynamic mixed-port kernel and phase-lag no-go & \StatusExactClosure{} & Dynamic \(3\times3\) bundle, determinant \(\Delta_\Pi D_\Pi\), dynamic product-family theorem, outgoing-port derivative identity, and first-order phase-lag/no-barrier theorem. \\
\addlinespace
204 & Resonance/linewidth tradeoff & \StatusExactClosure{} & Simple-pole Breit--Wigner normal form, dispersive/absorptive no-free-lunch theorem, low-loss bound, and linear survival window. \\
\addlinespace
205 & Concrete finite-throat primitive branch & Mixed: \StatusExactClosure{}, \StatusNumerical{} & Lowest N/N and D/N primitive branch, quartic pole polynomial, pole census, residue/linewidth figure, and first branch-level dynamic survival test. \\
\addlinespace
206 & Isotropic target surface and primitive compatibility & Mixed: \StatusExactClosure{}, \StatusNumerical{} & Exact compatibility equation between the conservative one-pole condition and outgoing normalization, compatibility branch, and finite dynamic survival windows. \\
\addlinespace
207 & Branch-packet compiler and actual-branch kill test & Mixed: \StatusExactClosure{}, \StatusOpen{} & Compiler from \((\Delta_{\rm norm},a_{P_0},b_{P_0})\) to lane prefactors, weak-axisymmetric law \(b_{P_0}=3a_{P_0}\), and transported ceiling in \((\Delta_{\rm norm},\Xi_1)\). \\
\addlinespace
208 & Microscopic \(\Xi_1\) compiler and first compensation surface & Mixed: \StatusExactClosure{}, \StatusOpen{} & First-order formulas for \(u_2^{(1)},u_4^{(1)},\Xi_1\), primitive logarithmic-slope compiler, mechanism sieve, and first surviving mixed-sector family. \\
\addlinespace
209 & Strict 5PN even-gate package & Mixed: \StatusExactClosure{}, \StatusReduced{} & Imports \(\Xi_{\rm load},K_1,H_{\rm even}\), proves \(\Xi_{\rm load}=\Xi_1\), and narrows the survivor to a pure outgoing-transfer corridor with the conservative bundle frozen. \\
\addlinespace
210 & Pure-transfer load factor and co-loading no-go & \StatusExactClosure{} & Exact factorization \(\Lambda=P/\Delta\), identity \(\Xi_1=2\delta\ln\Lambda\), outgoing-rigidity sieve, and co-loading no-go under simultaneous interference and hybridization rigidity. \\
\addlinespace
211 & Numerator/denominator split and dynamic-window test & Mixed: \StatusExactClosure{}, \StatusNumerical{} & Split \(\Xi_1=2(\pi_1-\delta_1)\), numerator-rigid and denominator-rigid survivors, and first wall-like dynamic-window audit. \\
\addlinespace
212 & Selected-branch numerator/denominator signature & \StatusExactClosure{} & Selected softening-depth factorization, classifier map, threshold \(\delta=8/9\), and denominator-like near-softening theorem. \\
\addlinespace
213 & Dynamic-window classifier and static-first theorem & Mixed: \StatusExactClosure{}, \StatusNumerical{} & Rigid-split share compiler, dynamic sign theorem, classifier threshold, and proof that the static \(\Xi_1\) budget remains the first kill condition. \\
\addlinespace
214 & Continuum-placement pullback & \StatusExactClosure{} & Pulls the selected-branch classifier back to \((\epsilon_\eta,\epsilon_W,\rho,Z_W,\delta_0,\Lambda)\), proves monotonicity in \(R_{\rm target}\), and keeps the static-first verdict on the physical placement map. \\
\addlinespace
215 & Known 5PN data injection & Mixed: \StatusNumerical{}, \StatusOpen{} & Injects numerically located Family-1 support/source data, shows the support/source side is safe by a large margin, and leaves the unresolved gate at static orbit-lock/coherent placement. \\
\addlinespace
216 & Rigid-mouth orbit-lock compiler & Mixed: \StatusExactClosure{}, \StatusOpen{} & Distinguishes mouth tracking from operator rigidity, shows \(\delta\ln R_{\rm tr}=0\Rightarrow\Xi_1=\delta\ln\mathfrak N_*\), and reads the static ceiling as an internal transfer/turbulence gate. \\
\addlinespace
217 & Direct branch-observable static gate & \StatusExactClosure{} & Finite chart in \((R_{\rm tr},R_{\rm target},\epsilon_\eta)\), two cancellations, rigid-mouth reduction, and two-observable kill test for \(R_{\rm target}\) and \(\epsilon_\eta\). \\
\addlinespace
218 & Rigid-mouth packet projectors & \StatusExactClosure{} & Involutive direct-packet map, projectors \(P_{\rm nt},P_\eta\), static-blind dressing line, and codimension-two orbit-lock theorem. \\
\addlinespace
219 & Microscopic dependent-plane projectors & \StatusExactClosure{} & Dependent-plane compiler, projectors \(P_{\rm nt}^{\rm dep},P_\eta^{\rm dep}\), equal-drift dressing ray, and static-only restoration gap. \\
\addlinespace
220 & Actual-branch dressing compiler & \StatusExactClosure{} & Finite static-blind curve, direct computation of \(q_\eta\), microscopic extractor \(q_\eta=2c_1-\kappa_U-\kappa_\eta\), and support-blindness theorem. \\
\addlinespace
221 & Physical transfer-shape compiler & \StatusExactClosure{} & Coherent identity \(R_{\rm target}\mathcal T^2=\Lambda_0(1-\epsilon_\eta)\), packet factorization, and three-gate theorem: tracking, transfer-shape, dressing. \\
\addlinespace
222 & Rigid-mouth physical normal form & \StatusExactClosure{} & Cartesian chart \((U,V)\), physical-to-microscopic compiler \((\Delta_T,\Delta_{K_\eta},\Delta_\mu)=(0,-V,U-V)\), and orbit-lock theorem \(U=V=0\). \\
\addlinespace
223 & Selected-branch loading ratio & \StatusExactClosure{} & Exact loading ratio \(\rho_\alpha=4/3\), required support excess \(\zeta_{\rm req}=1/3\), and selection of the symmetric lowest-twin support slice. \\
\addlinespace
224 & Primitive ranking on selected twin-support branch & \StatusExactClosure{} & Exact selected curve \(\epsilon_*=1-3\varrho/2\), threshold pair \((\varrho_{W\Lambda},\varrho_{U\Lambda})\), and primitive carrier ranking phase diagram. \\
\addlinespace
225 & Actual twin-support placement and coherent orbit-lock compiler & Mixed: \StatusExactClosure{}, \StatusOpen{} & Actual placement coordinate \(\varrho_{\rm phys}=2(1-\epsilon)/3\), support/orbit packet split, and final front-end branch-evaluation packet. \\
\bottomrule
\end{longtable}

\claimstatus{Part VII is mixed in status.  The static and dynamic one-port kernels, the projectors, the physical chart, and the selected support-ratio compilers are \StatusExactClosure{} inside the declared one-port, weak-axisymmetric, selected-branch, and rigid-mouth closures.  Numerical statements in Stages 205--206 and Stage 215 are \StatusNumerical{} insertions into exact formulas.  Any claim that the completed nonlinear moving-throat PDE actually realizes the winning static placement, dynamic window, rigid-mouth orbit-lock point, or selected twin-support point remains \StatusOpen{}.  No result in this part should be cited as a proof of a new same-charge long-range attraction.}

\section{Block contract and theorem-level output}
\label{sec:app-part07-block-contract}

The contract of Part VII is to separate three questions that are easy to conflate.  First, there is the \emph{kernel-class question}: does the mixed sector create a new long-range same-charge attractive law?  Stages 202--204 answer no within the one-port linear closure.  Second, there is the \emph{window question}: can a short-range static or resonant dynamic corridor survive the exact transported ceilings?  Stages 205--215 reduce that question to finite branch inequalities and show that the support/source side is not the active bottleneck in the currently injected data.  Third, there is the \emph{actual branch question}: does the PDE-selected rigid-mouth branch land on the orbit-lock and selected-support packets?  Stages 216--225 compile that question into direct observables.

The operational chain is
\begin{equation}
\label{eq:app-part07-operational-chain}
\begin{aligned}
\mathcal K_{\rm red}
&\longrightarrow
\mathfrak V_{\rm mix}(x,\omega)
\longrightarrow
(\Delta_{\rm norm},\Xi_1)
\\
&\longrightarrow
(R_{\rm tr},\mathcal T^2,\epsilon_\eta)
\longrightarrow
(U,V)
\\
&\longrightarrow
\text{selected twin-support placement packet}.
\end{aligned}
\end{equation}

\begin{theorem}[Part VII same-charge and rigid-mouth selected-branch theorem]
\label{thm:app-part07-main}
Inside the declared one-port wall/BdG/Maxwell/mixed closure, weak-axisymmetric grouped response closure, rigid-mouth actual-branch closure, and selected twin-support closure, the following statements hold.
\begin{enumerate}[leftmargin=2.2em]
  \item The static same-charge one-port bundle is a positive quadratic response on the admissible side \(\Omega_U^2>0\), \(\Delta>0\), \(D_0>0\).  It produces only product kernels built from the primitive source profiles.  For the carried source pair \(x^{-3}\) and \(e^{-2\kappa x}/x\), the only static product families are
  \begin{equation}
  \label{eq:app-part07-main-static-families}
  x^{-6},
  \qquad
  e^{-2\kappa x}x^{-4},
  \qquad
  e^{-4\kappa x}x^{-2}.
  \end{equation}
  \item The linear dynamic mixed bundle preserves the same spatial families.  A passive/outgoing correction at first order is pure phase lag and carries no real conservative barrier lowering off pole.
  \item Near a simple conservative pole, every relevant susceptibility has the local form
  \begin{equation}
  \label{eq:app-part07-main-line-shape}
  \chi_*(\omega)=\frac{A_*}{\delta-i\gamma_*},
  \qquad
  \delta=\omega-\omega_*,
  \end{equation}
  so
  \begin{equation}
  \label{eq:app-part07-main-line-ratio}
  \frac{|\Re\chi_*|}{|\Im\chi_*|}=\frac{|\delta|}{\gamma_*}.
  \end{equation}
  The maximum conservative gain occurs exactly when conservative and absorptive magnitudes are equal.
  \item On the explicit finite-throat primitive branch, the conservative one-pole condition and outgoing normalization condition are simultaneously compatible only on the exact surface
  \begin{equation}
  \label{eq:app-part07-main-compatibility-surface}
  \frac{N_0}{P_{0,{\rm target}}}
  =
  \frac{3(M+B_2+Z_2)^2}{B_4+Z_4}.
  \end{equation}
  The transported dynamic window is finite on that surface.
  \item The actual weak-axisymmetric branch test is the transported prefactor packet
  \begin{equation}
  \label{eq:app-part07-main-prefactor-packet}
  \bigl(\Delta_{\rm norm},a_{P_0},b_{P_0}\bigr),
  \qquad
  b_{P_0}=3a_{P_0},
  \qquad
  \Xi_1=\frac{P_1}{P_0}.
  \end{equation}
  The robust one-scalar ceiling is
  \begin{equation}
  \label{eq:app-part07-main-xi-ceiling}
  \frac{\Delta_{\rm norm}+T_{\rm quad}}{\hat m_0^{\,2}}
  \bigl(1+|\varepsilon\Xi_1|\bigr)
  \le P_{\rm crit},
  \qquad
  T_{\rm quad}:=\frac{54Gc_s^5}{5a^5c^5}.
  \end{equation}
  \item On the rigid-mouth actual branch, the physical packet is diagonal in the variables \((U,V)\).  The microscopic dependent correction is
  \begin{equation}
  \label{eq:app-part07-main-dependent-correction}
  (\Delta_T,\Delta_{K_\eta},\Delta_\mu)=(0,-V,U-V).
  \end{equation}
  Therefore the static-only correction changes only \(\mu_W\), while the post-static dressing correction is the equal \(K_\eta^{(\rm eff)}\)--\(\mu_W\) ray.
  \item Rigid-mouth orbit lock is the codimension-two Cartesian condition
  \begin{equation}
  \label{eq:app-part07-main-cartesian-lock}
  U=0,
  \qquad
  V=0,
  \end{equation}
  equivalently
  \begin{equation}
  \label{eq:app-part07-main-cartesian-lock-physical}
  \mathcal T^2=\mathcal T_{\rm ref}^2,
  \qquad
  \epsilon_\eta=\epsilon_{\eta,\rm ref}.
  \end{equation}
  \item The selected passive/outgoing support side lands exactly on the symmetric lowest-twin slice
  \begin{equation}
  \label{eq:app-part07-main-twin-slice}
  \rho_\alpha=\frac43,
  \qquad
  \zeta_{\rm req}=\frac13,
  \qquad
  C_{\rm mix}<\Pi_{\rm tr}<2C_{\rm mix}.
  \end{equation}
  \item The actual moving-throat front-end packet to evaluate is
  \begin{equation}
  \label{eq:app-part07-main-final-packet}
  \left(
  \epsilon,
  \varrho_{\rm phys},
  \sigma_{\rm phys},
  \text{ranking region},
  R_{\rm tr},
  R_{\rm target},
  \epsilon_\eta,
  d\ln R_{\rm tr},
  d\ln R_{\rm target},
  d\ln\epsilon_\eta,
  N_Q-1
  \right).
  \end{equation}
\end{enumerate}
Any failure of the actual branch to realize \cref{eq:app-part07-main-cartesian-lock,eq:app-part07-main-twin-slice,eq:app-part07-main-final-packet} is a branch-realization failure, not a failure of the algebraic compiler.
\end{theorem}

\section{One-port same-charge static kernel}
\label{sec:app-part07_same_charge_rigid_mouth:one-port-same-charge-static-kernel}

Stage 202 begins with the same isotropic one-port bundle already used by the grouped normalization chain.  After the stable BdG support mode is integrated out, the effective static wall stiffness is
\begin{equation}
\label{eq:app-part07-kstar}
K_*:=K-\frac{C^2}{\varpi^2}.
\end{equation}
The one-port Maxwell/mixed block is controlled by
\begin{equation}
\label{eq:app-part07-static-delta-qp}
\Delta:=\Omega_U^2\Omega_W^2-R^2,
\qquad
Q:=G_U^2\Omega_W^2+2G_UG_WR+G_W^2\Omega_U^2,
\qquad
P:=\Omega_U^2G_W+RG_U,
\end{equation}
with static wall operator
\begin{equation}
\label{eq:app-part07-static-d0}
D_0:=K_* - \frac{Q}{\Delta}.
\end{equation}
The same quantities control the outgoing prefactor:
\begin{equation}
\label{eq:app-part07-static-n0-p0}
N_0=\frac{P^2}{\Delta^2},
\qquad
P_0=\frac{N_0}{D_0}.
\end{equation}
Thus the static same-charge audit is not independent of the quadrupole normalization program; it probes the same reduced bundle.

The reduced static energy is
\begin{equation}
\label{eq:app-part07-static-energy}
V_{\rm stat}(q,U,W;r)
=
\frac12K_*q^2+\frac12\Omega_U^2U^2+\frac12\Omega_W^2W^2
-RUW-G_UqU-G_WqW-J_q(r)q-J_U(r)U-J_W(r)W.
\end{equation}
In matrix form,
\begin{equation}
\label{eq:app-part07-static-kred}
V_{\rm stat}=\frac12X^T\mathcal K_{\rm red}X-J^TX,
\qquad
X=(q,U,W)^T,
\end{equation}
where
\begin{equation}
\label{eq:app-part07-static-kred-matrix}
\mathcal K_{\rm red}
=
\begin{pmatrix}
K_* & -G_U & -G_W \\
-G_U & \Omega_U^2 & -R \\
-G_W & -R & \Omega_W^2
\end{pmatrix}.
\end{equation}
The determinant identity is
\begin{equation}
\label{eq:app-part07-static-det}
\det\mathcal K_{\rm red}=\Delta D_0.
\end{equation}
The admissible static side is therefore
\begin{equation}
\label{eq:app-part07-static-admissible}
\Omega_U^2>0,
\qquad
\Delta>0,
\qquad
D_0>0.
\end{equation}
On this side, minimization gives
\begin{equation}
\label{eq:app-part07-static-onshell}
X_*(r)=\mathcal K_{\rm red}^{-1}J(r),
\qquad
\delta V_{\rm mix}(r)=-\frac12J(r)^T\mathcal K_{\rm red}^{-1}J(r)\le0.
\end{equation}
So the one-port static response is attractive or neutral at quadratic order.  But its spatial family is inherited from the source products, not newly generated.

The wall--mixed cross susceptibility is especially important:
\begin{equation}
\label{eq:app-part07-static-chi-qw}
\chi_{qW}:=(\mathcal K_{\rm red}^{-1})_{qW}
=
\frac{P}{\Delta D_0}.
\end{equation}
With \(\Lambda:=P/\Delta\), one has
\begin{equation}
\label{eq:app-part07-static-load-identities}
N_0=\Lambda^2,
\qquad
P_0=\frac{\Lambda^2}{D_0},
\qquad
\chi_{qW}=\frac{\Lambda}{D_0},
\qquad
\chi_{qW}^2=\frac{P_0}{D_0}.
\end{equation}
So static mixed softening and outgoing quadrupole normalization are linked by the same load factor and denominator.

For the primitive source profiles
\begin{equation}
\label{eq:app-part07-primitive-sources}
\mathcal S_Q(x)=\frac1{x^3},
\qquad
\mathcal S_Y(x)=\frac{e^{-2\kappa x}}{x},
\end{equation}
and the source vector
\begin{equation}
\label{eq:app-part07-primitive-source-vector}
J(x)=
\begin{pmatrix}
\beta_Q\mathcal S_Q(x)\\
\beta_U\mathcal S_Y(x)\\
\beta_W\mathcal S_Y(x)
\end{pmatrix},
\end{equation}
the static correction is
\begin{equation}
\label{eq:app-part07-static-product-kernel}
\widetilde{\delta V}_{\rm mix}(x)
=-\frac12\left[
\frac{\mathcal C_6}{x^6}
+2\mathcal C_4\frac{e^{-2\kappa x}}{x^4}
+\mathcal C_2\frac{e^{-4\kappa x}}{x^2}
\right].
\end{equation}
The coefficients are susceptibility-weighted source products.  The key theorem is the family statement, not the coefficient details: no \(-1/x\), no \(-e^{-2\kappa x}/x\), and no longer-range static attraction appears inside this one-port static closure.

\begin{theorem}[Static one-port square-law suppression]
\label{thm:app-part07-static-square-law}
Inside the admissible one-port static closure, the mixed bundle creates no new long-range same-charge attractive law.  For the primitive source pair \(x^{-3}\) and \(e^{-2\kappa x}/x\), it produces only the three product families in \cref{eq:app-part07-static-product-kernel}.  Static same-charge survival is therefore a coefficient-engineering or near-softening question on already-short-range families, not a new spatial-kernel theorem.
\end{theorem}

\claimstatus{Stage 202 is \StatusExactClosure{} within the isotropic one-port static closure.  The absence of a new long-range family is exact for the stated source families and bundle assumptions.  The magnitude of the coefficients and their realized branch values remain branch data.}

\section{Linear dynamic mixed-bundle kernel}
\label{sec:app-part07_same_charge_rigid_mouth:linear-dynamic-mixed-bundle-kernel}

Stages 203--204 ask whether time dependence rescues what the static kernel does not provide.  The dynamic lift keeps the same three coordinates but replaces the static coefficients by
\begin{equation}
\label{eq:app-part07-dyn-coefficients}
K_B(\omega):=K-M\omega^2-\frac{C^2}{\varpi^2-\omega^2},
\qquad
A(\omega):=\Omega_U^2-\omega^2,
\qquad
W(\omega):=\Omega_W^2-\omega^2-\Pi(\omega).
\end{equation}
Then
\begin{equation}
\label{eq:app-part07-dyn-delta-q-d}
\Delta_\Pi(\omega):=A(\omega)W(\omega)-R^2,
\qquad
Q_\Pi(\omega):=G_U^2W(\omega)+2G_UG_WR+G_W^2A(\omega),
\end{equation}
and
\begin{equation}
\label{eq:app-part07-dyn-dpi}
D_\Pi(\omega):=K_B(\omega)-\frac{Q_\Pi(\omega)}{\Delta_\Pi(\omega)}.
\end{equation}
The dynamic stiffness matrix is
\begin{equation}
\label{eq:app-part07-dyn-kmatrix}
\mathcal K_{\rm dyn}(\omega)=
\begin{pmatrix}
K_B(\omega) & -G_U & -G_W \\
-G_U & A(\omega) & -R \\
-G_W & -R & W(\omega)
\end{pmatrix},
\end{equation}
with determinant
\begin{equation}
\label{eq:app-part07-dyn-det}
\det\mathcal K_{\rm dyn}(\omega)=\Delta_\Pi(\omega)D_\Pi(\omega).
\end{equation}
The complex response functional is
\begin{equation}
\label{eq:app-part07-dyn-vmix}
\mathfrak V_{\rm mix}(x,\omega)
=-\frac12J(x,\omega)^T\mathcal K_{\rm dyn}(\omega)^{-1}J(x,\omega).
\end{equation}
Its real part is the in-phase barrier reshaping; its imaginary part is the quadrature pumping/leakage channel.

For the same primitive source families as Stage 202, one gets the exact dynamic product-family theorem
\begin{equation}
\label{eq:app-part07-dyn-product-family}
\mathfrak V_{\rm mix}(x,\omega)
=-\frac12\left[
\frac{\mathcal C_6(\omega)}{x^6}
+2\mathcal C_4(\omega)\frac{e^{-2\kappa x}}{x^4}
+\mathcal C_2(\omega)\frac{e^{-4\kappa x}}{x^2}
\right].
\end{equation}
Thus linear monochromatic driving does not create a new radial family.  It promotes the same coefficients to complex susceptibilities.

The outgoing port enters only through the \(WW\) slot.  With \(e_W=(0,0,1)^T\),
\begin{equation}
\label{eq:app-part07-outgoing-derivative}
\partial_\Pi \mathcal K_{\rm dyn}^{-1}
=
\mathcal K_{\rm dyn}^{-1}e_We_W^T\mathcal K_{\rm dyn}^{-1},
\end{equation}
so
\begin{equation}
\label{eq:app-part07-vmix-pi-derivative}
\partial_\Pi\mathfrak V_{\rm mix}
=-\frac12\left[e_W^T\mathcal K_{\rm dyn}^{-1}J\right]^2.
\end{equation}
On the conservative off-pole branch the bracket is real.  Therefore, for a passive/outgoing port
\begin{equation}
\label{eq:app-part07-passive-pi}
\Pi(\omega)=i\Gamma(\omega)+O(\Gamma^2),
\qquad
\Gamma(\omega)\ge0,
\end{equation}
the first correction is
\begin{equation}
\label{eq:app-part07-phase-lag-correction}
\delta\mathfrak V_{\rm mix}^{(1)}(x,\omega)
=-\frac{i}{2}\Gamma(\omega)\mathcal T_J(\omega)^2,
\end{equation}
where \(\mathcal T_J=e_W^T\mathcal K_{\rm cons}^{-1}J\).  Consequently
\begin{equation}
\label{eq:app-part07-phase-lag-no-real}
\Re\delta\mathfrak V_{\rm mix}^{(1)}=0,
\qquad
\overline P_{\rm abs}^{(1)}(x,\omega)
=\frac{\omega\Gamma(\omega)}{2}\mathcal T_J(\omega)^2\ge0.
\end{equation}
This is the phase-lag no-go: passive/outgoing dressing gives absorption at first order, not conservative lowering.

Stage 204 then studies the only remaining dynamic escape hatch: resonant dispersive enhancement of the already-short-range families.  Near a simple conservative pole, write
\begin{equation}
\label{eq:app-part07-simple-pole-denominator}
F_\Pi(\omega)=F_0'(\omega_*)\delta-\Pi(\omega_*)Z_*+O(\delta^2,\Pi\delta,\Pi^2),
\qquad
\delta=\omega-\omega_*.
\end{equation}
For \(\Pi(\omega_*)=i\Gamma_*\), the reduced susceptibility has the universal form
\begin{equation}
\label{eq:app-part07-breit-wigner}
\chi_s(\omega)\approx\frac{A_*}{\delta-i\gamma_*},
\qquad
A_*:=\frac{\mathcal N_s(\omega_*)}{F_0'(\omega_*)},
\qquad
\gamma_*:=\frac{\Gamma_*Z_*}{|F_0'(\omega_*)|}.
\end{equation}
This gives
\begin{equation}
\label{eq:app-part07-re-im-ratio}
|\Re\chi_*|=\frac{|A_*|}{\gamma_*}\frac{r}{1+r^2},
\qquad
|\Im\chi_*|=\frac{|A_*|}{\gamma_*}\frac1{1+r^2},
\qquad
r:=\frac{|\delta|}{\gamma_*}.
\end{equation}
The exact maximum of the conservative part occurs at \(r=1\), where \(|\Re\chi_*|=|\Im\chi_*|\).  If one demands a low-loss window
\begin{equation}
\label{eq:app-part07-low-loss-window}
|\Im\chi_*|\le\eta|\Re\chi_*|,
\qquad
0<\eta\le1,
\end{equation}
then the best possible conservative magnitude is
\begin{equation}
\label{eq:app-part07-low-loss-bound}
\sup_{|\Im\chi_*|\le\eta|\Re\chi_*|}|\Re\chi_*|
=
\frac{|A_*|}{\gamma_*}\frac{\eta}{1+\eta^2}.
\end{equation}
For a spatial family \(S_j(x)\), a necessary low-loss survival condition is
\begin{equation}
\label{eq:app-part07-linear-survival-window}
\frac{|A_j|}{\gamma_*}\frac{\eta}{1+\eta^2}S_j(x)^2
\ge
2\Delta V_{\rm req}(x).
\end{equation}

\begin{theorem}[Linear dynamic no-free-lunch theorem]
\label{thm:app-part07-dynamic-no-free-lunch}
Inside the one-port linear dynamic closure, time dependence does not create a new same-charge spatial kernel.  Passive/outgoing dressing is pure phase lag at first order off pole.  Near a pole, useful conservative gain and absorptive loading obey the exact line-shape tradeoff \cref{eq:app-part07-re-im-ratio}.  Therefore any surviving linear dynamic corridor must be a high-quality, residue-to-linewidth enhancement of one of the already-known short-range families.
\end{theorem}

\claimstatus{Stages 203--204 are \StatusExactClosure{} for the dynamic bundle algebra, phase-lag identity, simple-pole normal form, and low-loss survival inequality.  Whether a completed branch supplies a useful pole with large enough residue-to-linewidth ratio remains \StatusOpen{}.}

\section{Actual-branch prefactor ceiling}
\label{sec:app-part07_same_charge_rigid_mouth:actual-branch-prefactor-ceiling}

Stages 205--207 turn the general dynamic survival inequality into a concrete transported branch test.  Stage 205 chooses the finite interval \(s\in[0,L]\), the lowest N/N zero mode
\begin{equation}
\label{eq:app-part07-nn-zero-mode}
u_0(s)=\frac1{\sqrt L},
\end{equation}
and the lowest D/N half-wave
\begin{equation}
\label{eq:app-part07-dn-half-wave}
f_0(s)=\sqrt{\frac2L}\sin\frac{\pi s}{2L}.
\end{equation}
The overlap is
\begin{equation}
\label{eq:app-part07-primitive-kappa}
\kappa=\int_0^Lu_0(s)f_0(s)\,ds=\frac{2\sqrt2}{\pi}.
\end{equation}
Thus
\begin{equation}
\label{eq:app-part07-primitive-couplings}
C=\kappa\lambda_B,
\qquad
G_U=\lambda_U,
\qquad
G_W=\kappa\lambda_W,
\qquad
R=\kappa\lambda_R.
\end{equation}
The conservative pole condition is a quartic in \(y=\omega^2\):
\begin{equation}
\label{eq:app-part07-quartic-pole}
\begin{aligned}
F(y)=&\left((K-My)(\varpi^2-y)-C^2\right)
\left((\Omega_U^2-y)(\Omega_W^2-y)-R^2\right)\\
&-(\varpi^2-y)
\left(G_U^2(\Omega_W^2-y)+2G_UG_WR+G_W^2(\Omega_U^2-y)\right).
\end{aligned}
\end{equation}
So the primitive finite-throat branch has four conservative poles in the generic admissible case.

For a simple wall-like pole with the outgoing quadrupole port
\begin{equation}
\label{eq:app-part07-outgoing-gamma5}
\Pi_{\rm out}(\omega)=i\Gamma_5\omega^5,
\qquad
\Gamma_5=\frac{a^5}{27c_s^5},
\end{equation}
the residue/linewidth figure for the pure quadrupolar same-charge family simplifies to
\begin{equation}
\label{eq:app-part07-residue-linewidth-primitive}
\mathcal R_{Q,*}
=
\frac{|A_{Q,*}|}{\gamma_*}
=
\frac{27c_s^5}{a^5\omega_*^5N(\omega_*)}.
\end{equation}
The derivative of the pole cancels; the figure is controlled by the pole frequency and outgoing transfer factor.

Stage 206 then imposes the isotropic one-pole target surface.  On the isotropic branch,
\begin{equation}
\label{eq:app-part07-d-coefficients-isotropic}
D_0=K-B_0-Z_0,
\qquad
D_2=-(M+B_2+Z_2),
\qquad
D_4=-(B_4+Z_4).
\end{equation}
The one-pole identity and outgoing-normalization condition solve for the same wall stiffness.  Their equality is exactly
\begin{equation}
\label{eq:app-part07-compatibility-equation}
\frac{N_0}{P_{0,{\rm target}}}
=
\frac{3(M+B_2+Z_2)^2}{B_4+Z_4}.
\end{equation}
Equivalently, the primitive branch induces the compatible target
\begin{equation}
\label{eq:app-part07-compatible-target}
P_{0,{\rm target,compat}}
=
\frac{N_0(B_4+Z_4)}{3(M+B_2+Z_2)^2}.
\end{equation}
This is the exact compatibility surface from the conservative one-pole side to the outgoing-normalization side.

Stage 207 transports the finite dynamic survival windows from the primitive compatibility family into the actual branch packet.  The final 5PN branch packet contains
\begin{equation}
\label{eq:app-part07-branch-packet-prefactor}
\Delta_{\rm branch}
=(a_2,b_2,a_4,b_4,a_{P_0},b_{P_0},\Delta_{\rm pole},\Delta_{\rm norm}),
\end{equation}
where
\begin{equation}
\label{eq:app-part07-delta-norm}
\Delta_{\rm norm}=\hat m_0^{\,2}\bar P_0-T_{\rm quad},
\qquad
T_{\rm quad}:=\frac{54Gc_s^5}{5a^5c^5}.
\end{equation}
Hence
\begin{equation}
\label{eq:app-part07-pbar-compiler}
\bar P_0=\frac{\Delta_{\rm norm}+T_{\rm quad}}{\hat m_0^{\,2}}.
\end{equation}
The grouped inverse map gives
\begin{equation}
\label{eq:app-part07-prefactor-lanes}
P_{20}=\bar P_0+4a_{P_0},
\qquad
P_{21}=\bar P_0-a_{P_0}+b_{P_0},
\qquad
P_{22}=\bar P_0-a_{P_0}-b_{P_0}.
\end{equation}
For a ceiling \(P_{\rm crit}\), the all-lane transported condition is
\begin{equation}
\label{eq:app-part07-max-prefactor-ceiling}
\max\{P_{20},P_{21},P_{22}\}\le P_{\rm crit}.
\end{equation}
On the weak-axisymmetric branch,
\begin{equation}
\label{eq:app-part07-lambda-signature}
\lambda_{20}=1,
\qquad
\lambda_{21}=\frac12,
\qquad
\lambda_{22}=-1,
\end{equation}
and
\begin{equation}
\label{eq:app-part07-xi-prefactor-form}
P_A=\bar P_0\bigl(1+\varepsilon\lambda_A\Xi_1\bigr),
\qquad
\Xi_1=\frac{P_1}{P_0}.
\end{equation}
Therefore
\begin{equation}
\label{eq:app-part07-ap0-bp0-xi}
a_{P_0}=\frac{\varepsilon\bar P_0\Xi_1}{4},
\qquad
b_{P_0}=\frac{3\varepsilon\bar P_0\Xi_1}{4},
\end{equation}
and the all-lane ceiling collapses to
\begin{equation}
\label{eq:app-part07-xi-prefactor-ceiling}
\bar P_0\bigl(1+|\varepsilon\Xi_1|\bigr)\le P_{\rm crit}.
\end{equation}
Substituting \cref{eq:app-part07-pbar-compiler} gives the actual-branch form already stated in \cref{eq:app-part07-main-xi-ceiling}.

\begin{theorem}[Transported actual-branch prefactor ceiling]
\label{thm:app-part07-actual-branch-ceiling}
Inside the transported primitive-family window closure, an actual weak-axisymmetric branch clears the robust all-lane ceiling only if
\begin{equation}
\label{eq:app-part07-thm-ceiling}
\frac{\Delta_{\rm norm}+T_{\rm quad}}{\hat m_0^{\,2}}
\bigl(1+|\varepsilon\Xi_1|\bigr)
\le P_{\rm crit}.
\end{equation}
If the branch is exactly calibrated, \(\Delta_{\rm norm}=0\), the same condition becomes a lower bound on the source-map factor,
\begin{equation}
\label{eq:app-part07-source-map-bound}
\hat m_0^{\,2}
\ge
\frac{T_{\rm quad}(1+|\varepsilon\Xi_1|)}{P_{\rm crit}}.
\end{equation}
\end{theorem}

\claimstatus{Stages 205--207 are exact for the primitive branch formulas, compatibility surface, and packet compiler.  Pole censuses, compatibility-point values, and numerical headroom budgets are \StatusNumerical{}.  The use of primitive-family ceilings as actual-branch sufficient conditions is a transported reduced test, not an unconditional theorem of the full PDE.}

\section{Strict 5PN even-gate package}
\label{sec:app-part07_same_charge_rigid_mouth:strict-5-pn-even-gate-package}

Stages 208--209 insert the weak-axisymmetric prefactor into the stricter first-order 5PN package.  At first order write
\begin{equation}
\label{eq:app-part07-first-order-d-n}
D_A=D_0+\varepsilon\lambda_A D_{01}+O(\varepsilon^2),
\qquad
N_A=N_0+\varepsilon\lambda_A N_{01}+O(\varepsilon^2).
\end{equation}
Then
\begin{equation}
\label{eq:app-part07-xi-load}
\Xi_{\rm load}
=
\frac{N_{01}}{N_0}-\frac{D_{01}}{D_0}
=
\frac{P_1}{P_0}
=
\Xi_1.
\end{equation}
The stricter imported first-order even gates are
\begin{equation}
\label{eq:app-part07-even-gates}
K_1:=D_{21}+\frac{D_{01}}{9},
\qquad
H_{\rm even}:=D_{41}-\frac23D_{21}-\frac{D_{01}}{27}.
\end{equation}
Thus the same scalar \(\Xi_1\) that controls the same-charge prefactor ceiling is only one member of the strict first-order package.  A branch must also pass the conservative even gates if it is to remain on the imported 5PN one-pole/hidden-even structure.

Stage 208 first imposes a weaker compensation surface that preserves the conservative grouped response on the explicit branch.  It then builds a primitive logarithmic-slope compiler from
\begin{equation}
\label{eq:app-part07-primitive-log-slopes}
(x_K,x_M,x_{\lambda_B},x_{\varpi},x_{\lambda_U},x_{\lambda_W},x_{\lambda_R},x_{\Omega_U},x_{\Omega_W})
\end{equation}
to
\begin{equation}
\label{eq:app-part07-primitive-compiler-output}
D_{01},
\qquad
D_{21},
\qquad
D_{41},
\qquad
N_{01},
\qquad
\Xi_1.
\end{equation}
The first sieve shows that wall-only and pure-BdG-only mechanisms do not supply the desired nontrivial same-charge load once the conservative compensation is enforced.  A mixed-sector family survives.

Stage 209 imports \cref{eq:app-part07-even-gates} and tightens the survivor.  The sharpest surviving subcorridor has
\begin{equation}
\label{eq:app-part07-pure-transfer-corridor-basic}
D_{01}=D_{21}=D_{41}=0,
\qquad
\Xi_1=\Xi_{\rm load}=\frac{N_{01}}{N_0}.
\end{equation}
In words: after the strict even gates and conservative-shape preservation are imposed on the noncanonical compatibility branch, the remaining same-charge effect is not generic mixed anisotropy.  It is pure outgoing-transfer enhancement with the conservative one-pole bundle frozen at first order.

\begin{theorem}[Strict first-order package and pure-transfer survivor]
\label{thm:app-part07-strict-even-package}
Inside the imported first-order 5PN package, the same-charge load scalar is exactly \(\Xi_{\rm load}=P_1/P_0\).  Passing the strict even gates requires controlling \(K_1\) and \(H_{\rm even}\) in addition to \(\Xi_{\rm load}\).  On the explicit compatibility branch with Stage-208 conservative-shape preservation, the sharp surviving subcorridor is the pure-transfer corridor \cref{eq:app-part07-pure-transfer-corridor-basic}.
\end{theorem}

\claimstatus{Stages 208--209 are \StatusExactClosure{} for the first-order compilers and even-gate identities.  The mechanism sieve is exact within the primitive slope space, while the actual PDE realization of the pure-transfer corridor is \StatusOpen{}.}

\section{Pure-transfer corridor and co-loading no-go}
\label{sec:app-part07_same_charge_rigid_mouth:pure-transfer-corridor-and-co-loading-no-go}

Stages 210--211 analyze the pure-transfer survivor.  Define the outgoing-load factor
\begin{equation}
\label{eq:app-part07-load-factor-lambda}
\Lambda:=\frac{P}{\Delta}.
\end{equation}
In the pure-transfer corridor,
\begin{equation}
\label{eq:app-part07-xi-lambda}
\Xi_1=\frac{N_{01}}{N_0}=2\frac{P_{01}}{P}-2\frac{\Delta_{01}}{\Delta}=2\delta\ln\Lambda.
\end{equation}
The one-port factorization is
\begin{equation}
\label{eq:app-part07-lambda-factorization}
\Lambda
=
\frac{G_W}{\Omega_W^2}\frac{1+I}{1-H},
\qquad
I:=\frac{RG_U}{\Omega_U^2G_W},
\qquad
H:=\frac{R^2}{\Omega_U^2\Omega_W^2}.
\end{equation}
The three pieces are: the raw mixed leg \(G_W/\Omega_W^2\), the interference ratio \(I\), and the hybridization/blocking ratio \(H\).

The outgoing-rigidity sieve asks which of those pieces can remain fixed while \(\Xi_1\) survives.  If both interference and hybridization are rigid,
\begin{equation}
\label{eq:app-part07-co-loading-no-go-hypothesis}
\delta\ln I=0,
\qquad
\delta\ln H=0,
\end{equation}
then the simultaneous pure-transfer and conservative constraints leave only the trivial pure-transfer drift on the concrete branch.  This is the first co-loading no-go: the same-charge effect cannot be reduced to a naive mixed-leg slogan under both rigidities.

Stage 211 then splits the pure-transfer slope into numerator and denominator pieces:
\begin{equation}
\label{eq:app-part07-numerator-denominator-split}
\Xi_1=2(\pi_1-\delta_1),
\qquad
\pi_1:=\frac{P_{01}}{P},
\qquad
\delta_1:=\frac{\Delta_{01}}{\Delta}.
\end{equation}
The exact rigid-subcorridor theorem is:
\begin{enumerate}[leftmargin=2.2em]
  \item pure-transfer plus \(\pi_1=0\) leaves a one-dimensional numerator-rigid survivor;
  \item pure-transfer plus \(\delta_1=0\) leaves a one-dimensional denominator-rigid survivor;
  \item imposing both \(\pi_1=0\) and \(\delta_1=0\) kills the corridor.
\end{enumerate}
Both one-dimensional survivors pass the first wall-like dynamic-window audit on the explicit compatibility point.  The first active ceiling remains the transported static \(|\varepsilon\Xi_1|\) budget, not the dynamic wall-like window.

\begin{theorem}[Pure-transfer and rigidity sieve]
\label{thm:app-part07-pure-transfer-sieve}
Inside the strict pure-transfer corridor, the same-charge scalar is the logarithmic outgoing-load slope \(\Xi_1=2\delta\ln\Lambda\).  The factorization \cref{eq:app-part07-lambda-factorization} separates raw mixed-leg, interference, and hybridization contributions.  Simultaneous interference and hybridization rigidity kills the nontrivial pure-transfer drift on the concrete compatibility branch, while either numerator-rigid or denominator-rigid splitting alone leaves a one-dimensional survivor.
\end{theorem}

\claimstatus{Stages 210--211 are \StatusExactClosure{} for the pure-transfer identities and rigidity splits.  The branch-specific dynamic-window survival statements are \StatusNumerical{} insertions into exact ceiling formulas.}

\section{Selected-branch dynamic-window compiler}
\label{sec:app-part07_same_charge_rigid_mouth:selected-branch-dynamic-window-compiler}

Stages 212--215 replace the free numerator/denominator split by the selected-branch softening-depth language and then pull the result back to physical continuum coordinates.

The selected normalization product is
\begin{equation}
\label{eq:app-part07-selected-nminus}
N_-(x)=\frac{\beta_0s_-(x)^2}{\kappa_0^2(A-x)},
\qquad
0\le x<A,
\end{equation}
with selected overlap
\begin{equation}
\label{eq:app-part07-selected-sminus}
s_-(x)=
\frac{\left[\kappa_0^2(x+\Delta K_{\rm ax})+\kappa_1^2x\right]^2}
{\kappa_0^2(x+\Delta K_{\rm ax})^2+\kappa_1^2x^2}.
\end{equation}
With dimensionless variables
\begin{equation}
\label{eq:app-part07-xi-delta-selected}
\xi:=\frac{x}{A},
\qquad
\delta:=\frac{\Delta K_{\rm ax}}{A},
\end{equation}
Stage 212 defines a log-slope classifier \(\mathcal R_{ND}\) that compares the numerator-like and denominator-like content of the selected branch.  The exact threshold
\begin{equation}
\label{eq:app-part07-delta-eight-ninths}
\delta=\frac89
\end{equation}
separates always-denominator selected branches from branches that begin numerator-dominant.  The selected branch is never literally either rigid Stage-211 subcorridor; it is an exact co-loading product.  However it is always denominator-like near softening, and it is denominator-like on the entire stable branch for \(\delta\ge8/9\).

Stage 213 uses the classifier to mix the carried numerator-rigid and denominator-rigid dynamic data.  The share weights are
\begin{equation}
\label{eq:app-part07-share-weights}
w_{\rm num}=\frac{\mathcal R_{ND}}{1+\mathcal R_{ND}},
\qquad
w_{\rm den}=\frac{1}{1+\mathcal R_{ND}}.
\end{equation}
The dynamic sign theorem says the upper wall-like pole always worsens, while the lower wall-like pole flips sign only at one classifier threshold.  Most importantly, the selected-branch dynamic ceilings are everywhere weaker than the universal transported static ceilings.  Therefore the first kill condition remains the static \(|\varepsilon\Xi_1|\) budget.

Stage 214 pulls \(\mathcal R_{ND}\) back through the continuum placement map.  The physical classifier is
\begin{equation}
\label{eq:app-part07-physical-classifier}
\mathcal R_{\rm phys}(\delta,R_{\rm target})
:=
\mathcal R_{ND}\bigl(\xi_{\rm req}(\delta,R_{\rm target}),\delta\bigr),
\end{equation}
where \(\xi_{\rm req}\) is the stable solution of
\begin{equation}
\label{eq:app-part07-f-equals-rtarget}
F(\xi,\delta)=R_{\rm target}.
\end{equation}
The monotonicity theorem is
\begin{equation}
\label{eq:app-part07-rphys-monotonic}
\partial_{R_{\rm target}}\mathcal R_{\rm phys}<0.
\end{equation}
So larger normalization demand pushes the physical selected branch in the denominator-like and dynamically safer direction.  The static-first theorem survives the pullback.

Stage 215 injects the known 5PN Family-1 support/source data.  The selected passive/outgoing branch requires
\begin{equation}
\label{eq:app-part07-zeta-req-rho-alpha}
\zeta_{\rm req}=\frac13,
\qquad
\rho_{\alpha}^{\rm req}=\frac43.
\end{equation}
The numerically located Family-1 support/source branch returns \(\zeta_{\rm phys}\approx2.4675\) for the carried wall-depth extractions, so the support/source side overshoots the canonical demand by a large margin.  This does not prove same-charge success; it says the first active unresolved bottleneck stays at the actual PDE-selected static orbit-lock/coherent placement point.

\begin{theorem}[Selected-branch static-first theorem]
\label{thm:app-part07-selected-static-first}
Inside the selected-branch dynamic-window compiler, the wall-like dynamic window does not become the first same-charge kill condition.  After continuum pullback and known support/source injection, the support/source side is non-bottlenecked in the carried data.  The live unresolved test remains the actual static orbit-lock/coherent placement packet.
\end{theorem}

\claimstatus{Stages 212--214 are \StatusExactClosure{} for the selected-branch classifier, share compiler, monotonic pullback, and static-first theorem.  Stage 215 is \StatusNumerical{} for the injected Family-1 support/source roots and margins.  The actual PDE-selected static placement remains \StatusOpen{}.}

\section{Rigid-mouth orbit-lock compiler}
\label{sec:app-part07_same_charge_rigid_mouth:rigid-mouth-orbit-lock-compiler}

Stages 216--219 translate the surviving static placement problem into rigid-mouth packet language.  The first firewall is that mouth rigidity and operator rigidity are not the same statement.  Rigid-mouth or tracking lock is represented by
\begin{equation}
\label{eq:app-part07-tracking-lock}
\delta\ln R_{\rm tr}=0,
\end{equation}
whereas \(D_{01}=0\) is the stronger algebraic statement that the effective grouped operator does not drift at first order.

The direct finite branch observables are
\begin{equation}
\label{eq:app-part07-direct-observables}
(R_{\rm tr},R_{\rm target},\epsilon_\eta).
\end{equation}
Relative to a reference branch, the finite quotient coordinates include
\begin{equation}
\label{eq:app-part07-qtr-qnt-qeta-finite}
q_{\rm tr}=-C_*\ln\frac{R_{\rm tr}}{R_{\rm tr,ref}},
\end{equation}
\begin{equation}
\label{eq:app-part07-qnt-finite}
q_{\rm nt}=B_*\ln\frac{R_{\rm tr}}{R_{\rm tr,ref}}
+
\ln\frac{1-\epsilon_\eta}{1-\epsilon_{\eta,\rm ref}}
-
\ln\frac{R_{\rm target}}{R_{\rm target,ref}},
\end{equation}
\begin{equation}
\label{eq:app-part07-qeta-finite}
q_\eta=\ln\frac{\epsilon_\eta}{\epsilon_{\eta,\rm ref}}.
\end{equation}
At first order, Stage 217 gives the direct compiler
\begin{equation}
\label{eq:app-part07-direct-first-order}
\Theta_1=\delta\ln R_{\rm tr},
\qquad
\Xi_1=-\delta\ln R_{\rm target}-c_\eta\delta\ln\epsilon_\eta,
\qquad
\mathcal R_1=\delta\ln R_{\rm target},
\end{equation}
with
\begin{equation}
\label{eq:app-part07-ceta}
c_\eta:=\frac{\epsilon_{\eta,*}}{1-\epsilon_{\eta,*}}.
\end{equation}
On the rigid-mouth slice, \(q_{\rm tr}=0\), so the static scalar is
\begin{equation}
\label{eq:app-part07-rigid-static-scalar}
\Xi_1=-R_1-c_\eta E_1,
\qquad
R_1:=\delta\ln R_{\rm target},
\qquad
E_1:=\delta\ln\epsilon_\eta.
\end{equation}
The two-observable static gate is therefore a condition on \((R_1,E_1)\), not on support/source amplitude.

Stage 218 makes the packet structure exact.  Define
\begin{equation}
\label{eq:app-part07-rigid-mouth-direct-vector}
\mathbf x_{\rm rm}:=\begin{pmatrix}R_1\\E_1\end{pmatrix},
\qquad
\mathbf q_{\rm rm}:=\begin{pmatrix}q_{\rm nt}\\q_\eta\end{pmatrix}
=\begin{pmatrix}\Xi_1\\E_1\end{pmatrix}.
\end{equation}
Then
\begin{equation}
\label{eq:app-part07-mrm}
\mathbf q_{\rm rm}=M_{\rm rm}\mathbf x_{\rm rm},
\qquad
M_{\rm rm}:=\begin{pmatrix}-1&-c_\eta\\0&1\end{pmatrix},
\qquad
M_{\rm rm}^2=I_2.
\end{equation}
The direct-space projectors are
\begin{equation}
\label{eq:app-part07-direct-projectors}
P_{\rm nt}=\begin{pmatrix}1&c_\eta\\0&0\end{pmatrix},
\qquad
P_\eta=\begin{pmatrix}0&-c_\eta\\0&1\end{pmatrix}.
\end{equation}
The static strip \(\Xi_1=0\) is codimension one.  Full rigid-mouth orbit lock is codimension two:
\begin{equation}
\label{eq:app-part07-rm-codim-two}
q_{\rm nt}=0,
\qquad
q_\eta=0
\quad\Longleftrightarrow\quad
R_1=0,
\qquad
E_1=0.
\end{equation}
The static-blind dressing line is
\begin{equation}
\label{eq:app-part07-static-blind-line}
\Xi_1=0
\quad\Longleftrightarrow\quad
\mathbf x_{\rm rm}=\begin{pmatrix}-c_\eta q_\eta\\q_\eta\end{pmatrix}.
\end{equation}
Thus clearing the first static same-charge ceiling need not make the true orbit-failure amplitude small.

Stage 219 identifies the microscopic carrier.  On the rigid-mouth slice the dependent correction is
\begin{equation}
\label{eq:app-part07-dependent-correction-q}
\begin{pmatrix}\Delta_T^{(q)}\\\Delta_{K_\eta}^{(q)}\\\Delta_\mu^{(q)}\end{pmatrix}
=
\begin{pmatrix}0\\-q_\eta\\q_{\rm nt}-q_\eta\end{pmatrix}.
\end{equation}
The static-blind line \(q_{\rm nt}=0\) maps to the equal-drift ray
\begin{equation}
\label{eq:app-part07-equal-drift-ray}
\begin{pmatrix}\Delta_T\\\Delta_{K_\eta}\\\Delta_\mu\end{pmatrix}
=
-q_\eta\begin{pmatrix}0\\1\\1\end{pmatrix}.
\end{equation}
So after the static strip is cleared, the remaining orbit defect is not generic throat motion.  It is one scalar equal-drift \(K_\eta^{(\rm eff)}\)--\(\mu_W\) dressing amplitude.

\begin{theorem}[Rigid-mouth packet and dependent-plane theorem]
\label{thm:app-part07-rigid-mouth-packet}
On the rigid-mouth slice, the first static same-charge gate tests only \(q_{\rm nt}=\Xi_1\).  The dressing coordinate \(q_\eta\) is invisible to that gate and maps microscopically to the equal-drift ray \(-q_\eta(0,1,1)^T\) in the dependent triple.  Therefore the static gate is necessary but not sufficient for orbit lock.
\end{theorem}

\claimstatus{Stages 216--219 are \StatusExactClosure{} for the rigid-mouth packet compilers and projectors.  The physical assumption that the actual branch is rigid-mouth/track-locked, and the realized value of \(q_\eta\), remain \StatusOpen{} branch data.}

\section{Physical logarithmic chart \texorpdfstring{\((U,V)\)}{(U,V)}}
\label{sec:app-part07_same_charge_rigid_mouth:physical-logarithmic-chart-u-v}

Stages 220--222 diagonalize the rigid-mouth packet in physical branch variables.  Stage 220 begins from the finite rigid-mouth packet
\begin{equation}
\label{eq:app-part07-rm-finite-packet}
q_{\rm nt}=\ln\frac{1-\epsilon_\eta}{1-\epsilon_{\eta,\rm ref}}
-
\ln\frac{R_{\rm target}}{R_{\rm target,ref}},
\qquad
q_\eta=\ln\frac{\epsilon_\eta}{\epsilon_{\eta,\rm ref}}.
\end{equation}
The first static gate \(q_{\rm nt}=0\) is the finite curve
\begin{equation}
\label{eq:app-part07-static-blind-curve}
\frac{R_{\rm target}}{R_{\rm target,ref}}
=
\frac{1-\epsilon_\eta}{1-\epsilon_{\eta,\rm ref}}.
\end{equation}
Parameterizing by \(q_\eta\),
\begin{equation}
\label{eq:app-part07-static-blind-curve-qeta}
\epsilon_\eta=\epsilon_{\eta,\rm ref}e^{q_\eta},
\qquad
\frac{R_{\rm target}}{R_{\rm target,ref}}
=
\frac{1-\epsilon_{\eta,\rm ref}e^{q_\eta}}{1-\epsilon_{\eta,\rm ref}}.
\end{equation}
Conversely, once the static gate is cleared,
\begin{equation}
\label{eq:app-part07-qeta-from-rtarget}
q_\eta=
\ln\left(
\frac{1-(1-\epsilon_{\eta,\rm ref})R_{\rm target}/R_{\rm target,ref}}
{\epsilon_{\eta,\rm ref}}
\right).
\end{equation}
The microscopic extractor is
\begin{equation}
\label{eq:app-part07-qeta-micro-finite}
q_\eta
=2\ln\frac{c_{\eta U}}{c_{\eta U,\rm ref}}
-
\ln\frac{K_U}{K_{U,\rm ref}}
-
\ln\frac{K_\eta^{(\rm eff)}}{K_{\eta,\rm ref}^{(\rm eff)}}.
\end{equation}
At first order,
\begin{equation}
\label{eq:app-part07-qeta-first-order}
q_\eta=2c_1-\kappa_U-\kappa_\eta.
\end{equation}
The dressing coordinate is support-blind:
\begin{equation}
\label{eq:app-part07-qeta-support-blind}
\partial_\zeta q_\eta=0,
\qquad
\partial_{M_{\rm tr}}q_\eta=0.
\end{equation}

Stage 221 then uses the coherent transfer-shape identity
\begin{equation}
\label{eq:app-part07-transfer-shape-identity}
R_{\rm target}\mathcal T^2=\Lambda_0(1-\epsilon_\eta).
\end{equation}
This converts the same-charge packet into physical variables.  Relative to a reference branch,
\begin{equation}
\label{eq:app-part07-qnt-transfer-shape}
q_{\rm nt}+\frac{B_*}{C_*}q_{\rm tr}
=
\ln\frac{\mathcal T^2}{\mathcal T_{\rm ref}^2}.
\end{equation}
On the rigid-mouth slice, \(q_{\rm tr}=0\), so
\begin{equation}
\label{eq:app-part07-rm-qnt-qeta-transfer}
q_{\rm nt}=\ln\frac{\mathcal T^2}{\mathcal T_{\rm ref}^2},
\qquad
q_\eta=\ln\frac{\epsilon_\eta}{\epsilon_{\eta,\rm ref}}.
\end{equation}
Thus the first static same-charge gate is exactly transfer-shape rigidity:
\begin{equation}
\label{eq:app-part07-transfer-shape-rigidity}
q_{\rm nt}=0
\quad\Longleftrightarrow\quad
\mathcal T^2=\mathcal T_{\rm ref}^2.
\end{equation}
The post-static gate is dressing rigidity:
\begin{equation}
\label{eq:app-part07-dressing-rigidity}
q_\eta=0
\quad\Longleftrightarrow\quad
\epsilon_\eta=\epsilon_{\eta,\rm ref}.
\end{equation}

Stage 222 packages this in the Cartesian variables
\begin{equation}
\label{eq:app-part07-uv-def}
U:=\ln\frac{\mathcal T^2}{\mathcal T_{\rm ref}^2},
\qquad
V:=\ln\frac{\epsilon_\eta}{\epsilon_{\eta,\rm ref}}.
\end{equation}
Then
\begin{equation}
\label{eq:app-part07-rm-packet-uv}
q_{\rm nt}=U,
\qquad
q_\eta=V.
\end{equation}
The target ratio factorizes as
\begin{equation}
\label{eq:app-part07-rtarget-uv}
\frac{R_{\rm target}}{R_{\rm target,ref}}
=
\frac{1-\epsilon_{\eta,\rm ref}e^V}{1-\epsilon_{\eta,\rm ref}}e^{-U}.
\end{equation}
The dependent microscopic correction is
\begin{equation}
\label{eq:app-part07-uv-dependent-correction}
\begin{pmatrix}\Delta_T\\\Delta_{K_\eta}\\\Delta_\mu\end{pmatrix}
=
\begin{pmatrix}0\\-V\\U-V\end{pmatrix}.
\end{equation}
A left inverse reads
\begin{equation}
\label{eq:app-part07-uv-left-inverse}
U=\Delta_\mu-\Delta_{K_\eta},
\qquad
V=-\Delta_{K_\eta}.
\end{equation}
The pure transfer-shape axis has image
\begin{equation}
\label{eq:app-part07-pure-transfer-image}
(U,0)
\longmapsto
(0,0,U),
\end{equation}
so clearing the static defect changes only \(\mu_W\).  The pure dressing axis has image
\begin{equation}
\label{eq:app-part07-pure-dressing-image}
(0,V)
\longmapsto
-V(0,1,1),
\end{equation}
the equal \(K_\eta^{(\rm eff)}\)--\(\mu_W\) ray.  Therefore full orbit lock is
\begin{equation}
\label{eq:app-part07-uv-orbit-lock}
U=0,
\qquad
V=0.
\end{equation}
At first order,
\begin{equation}
\label{eq:app-part07-uv-first-order-dependent}
\begin{pmatrix}\Delta_T\\\Delta_{K_\eta}\\\Delta_\mu\end{pmatrix}^{(1)}
=
\begin{pmatrix}
0\\
-\delta\ln\epsilon_\eta\\
\delta\ln\mathcal T^2-\delta\ln\epsilon_\eta
\end{pmatrix}.
\end{equation}

\begin{theorem}[Rigid-mouth Cartesian orbit-lock theorem]
\label{thm:app-part07-cartesian-orbit-lock}
On the rigid-mouth actual branch, the physical logarithmic coordinates \((U,V)\) are the finite quotient packet.  The first static gate is \(U=0\).  The post-static dressing gate is \(V=0\).  The full orbit-lock point is the Cartesian codimension-two condition \((U,V)=(0,0)\), equivalently \(\mathcal T^2=\mathcal T_{\rm ref}^2\) and \(\epsilon_\eta=\epsilon_{\eta,\rm ref}\).
\end{theorem}

\claimstatus{Stages 220--222 are \StatusExactClosure{} for the finite static-blind curve, physical transfer-shape factorization, support-blindness, Cartesian chart, dependent correction, and orbit-lock theorem.  Whether the actual branch has \(U=V=0\) remains \StatusOpen{}.}

\section{Selected twin-support branch and placement}
\label{sec:app-part07_same_charge_rigid_mouth:selected-twin-support-branch-and-placement}

Stages 223--225 close the selected support side and attach it to the actual branch packet.  Stage 223 begins with the selected product identities
\begin{equation}
\label{eq:app-part07-selected-product-identities}
\Pi_{\rm tr}=\frac{N_Q^{(\rm target)}}{\beta_0}\alpha_{\rm req},
\qquad
C_{\rm mix}=\frac{N_Q^{(\rm target)}}{\beta_0}\alpha_{\rm mix}.
\end{equation}
Therefore
\begin{equation}
\label{eq:app-part07-rho-alpha-ratio}
\frac{\Pi_{\rm tr}}{C_{\rm mix}}
=\frac{\alpha_{\rm req}}{\alpha_{\rm mix}}
=:\rho_\alpha.
\end{equation}
For the minimal isotropic conservative quadrupole precursor
\begin{equation}
\label{eq:app-part07-minimal-quadrupole-precursor}
Y_Q^{\rm cons}(\omega)
=c_0+\frac{c_1}{1-\omega^2/\Omega_Q^2},
\qquad
c_0=\frac34,
\qquad
c_1=\frac14,
\qquad
\Omega_Q=\frac{3c_s}{2a},
\end{equation}
the contact/pole inverse compiler gives
\begin{equation}
\label{eq:app-part07-rho-alpha-zeta-req}
\rho_\alpha=\frac1{c_0}=\frac43,
\qquad
\zeta_{\rm req}:=\rho_\alpha-1=\frac13.
\end{equation}
Therefore
\begin{equation}
\label{eq:app-part07-pi-tr-four-thirds}
\Pi_{\rm tr}=\frac43 C_{\rm mix}.
\end{equation}
With
\begin{equation}
\label{eq:app-part07-varrho-selector}
\varrho:=\frac{\pi^2\Pi_{\rm tr}}{16\Lambda},
\qquad
C_{\rm mix}=\frac{8\Lambda(1-\epsilon_*)}{\pi^2},
\end{equation}
this becomes
\begin{equation}
\label{eq:app-part07-varrho-selected}
\varrho=\frac{2(1-\epsilon_*)}{3},
\qquad
S_{\rm req}=\frac43.
\end{equation}
Thus the selected branch lies strictly in the symmetric lowest-twin window:
\begin{equation}
\label{eq:app-part07-twin-window}
C_{\rm mix}<\Pi_{\rm tr}<2C_{\rm mix}.
\end{equation}

Stage 224 turns this into a primitive ranking phase diagram.  The selected curve is
\begin{equation}
\label{eq:app-part07-selected-curve}
\epsilon_* = 1-\frac{3\varrho}{2},
\qquad
\sigma=\frac{4}{3\varrho}-2,
\qquad
0<\varrho<\frac23.
\end{equation}
For the constructive coherent parameter \(0<\beta<2/11\), only two crossover thresholds survive:
\begin{equation}
\label{eq:app-part07-ranking-thresholds}
\varrho_{W\Lambda}
=
\frac{2(1+\beta^2)}{3(2+\beta^2)},
\qquad
\varrho_{U\Lambda}
=
\frac{2(1+\beta^2)}{3(1+\beta+\beta^2)}.
\end{equation}
They obey
\begin{equation}
\label{eq:app-part07-threshold-order}
0<\varrho_{W\Lambda}<\varrho_{U\Lambda}<\frac23.
\end{equation}
The ranking theorem is
\begin{equation}
\label{eq:app-part07-ranking-region-i}
0<\varrho<\varrho_{W\Lambda}
\quad\Longrightarrow\quad
q_\chi>q_Z>q_\Lambda>q_W>|q_U|,
\end{equation}
\begin{equation}
\label{eq:app-part07-ranking-region-ii}
\varrho_{W\Lambda}<\varrho<\varrho_{U\Lambda}
\quad\Longrightarrow\quad
q_\chi>q_Z>q_W>q_\Lambda>|q_U|,
\end{equation}
\begin{equation}
\label{eq:app-part07-ranking-region-iii}
\varrho_{U\Lambda}<\varrho<\frac23
\quad\Longrightarrow\quad
q_\chi>q_Z>q_W>|q_U|>q_\Lambda.
\end{equation}
So the selected branch is dominated by interference/overlap/wall-blocking/outgoing-scale corrections, not by a large split-\(U\) correction.

Stage 225 finally places the actual coherent branch on this curve.  The selected physical support coordinate is
\begin{equation}
\label{eq:app-part07-epsilon-physical}
\epsilon
=
\epsilon_W\left(1-\frac{2}{11}\frac{\delta_U}{1+\delta_U}\right).
\end{equation}
Therefore
\begin{equation}
\label{eq:app-part07-varrho-phys-sigma-phys}
\varrho_{\rm phys}=\frac23(1-\epsilon),
\qquad
\sigma_{\rm phys}=\frac{2\epsilon}{1-\epsilon}
=\frac{4}{3\varrho_{\rm phys}}-2.
\end{equation}
The coherent placement observables are
\begin{equation}
\label{eq:app-part07-rtr-coherent}
R_{\rm tr}
=
\frac{1+\chi_0/(1+\delta_U)}{1+\chi_0},
\end{equation}
\begin{equation}
\label{eq:app-part07-rtarget-coherent}
R_{\rm target}
=
\frac{\Lambda(1-\epsilon_\eta)(1-\epsilon)^2}{Z_W(1+\chi_0)^2},
\end{equation}
together with \(\epsilon_\eta\).  The finite coherent orbit packet in physical placement variables is
\begin{equation}
\label{eq:app-part07-stage225-qtr}
q_{\rm tr}
=(1+\delta_{U,*})\ln\frac{\chi_0}{\chi_{0,\rm ref}}
+(1+\chi_{0,*})\ln\frac{\delta_U}{\delta_{U,\rm ref}},
\end{equation}
\begin{equation}
\label{eq:app-part07-stage225-qnt}
q_{\rm nt}
=
\ln\frac{Z_W}{Z_{W,\rm ref}}
-
\ln\frac{\Lambda}{\Lambda_{\rm ref}}
+E_*\ln\frac{\epsilon_W}{\epsilon_{W,\rm ref}}
-F_*\ln\frac{\delta_U}{\delta_{U,\rm ref}},
\end{equation}
\begin{equation}
\label{eq:app-part07-stage225-qeta}
q_\eta=\ln\frac{\epsilon_\eta}{\epsilon_{\eta,\rm ref}}.
\end{equation}
The infinitesimal orbit-lock conditions are equivalently
\begin{equation}
\label{eq:app-part07-stage225-orbit-lock-inf}
d\ln R_{\rm tr}=0,
\qquad
d\ln R_{\rm target}=0,
\qquad
d\ln\epsilon_\eta=0.
\end{equation}
The outgoing finish line is separate:
\begin{equation}
\label{eq:app-part07-stage225-nq}
N_Q=1,
\qquad
\text{or equivalently }\chi_Q=1\text{ on the natural source-map branch.}
\end{equation}
The practical front-end packet to be returned by the completed PDE is \cref{eq:app-part07-main-final-packet}; if the support lane is tracked separately, the support packet is
\begin{equation}
\label{eq:app-part07-stage225-support-packet}
\left(\zeta,
M_{\rm mix},
S(\zeta;\epsilon),
M_{\rm tr}
\right),
\qquad
S(\zeta;\epsilon)=1+\frac{\zeta(1-\epsilon)}{1-\zeta\epsilon}.
\end{equation}
The support variable \(\zeta\) affects the support baseline but not the orbit packet:
\begin{equation}
\label{eq:app-part07-zeta-orbit-blind}
\partial_\zeta R_{\rm tr}
=
\partial_\zeta R_{\rm target}
=
\partial_\zeta\epsilon_\eta
=
\partial_\zeta q_{\rm tr}
=
\partial_\zeta q_{\rm nt}
=
\partial_\zeta q_\eta
=0.
\end{equation}

\begin{theorem}[Selected twin-support placement and orbit-packet split]
\label{thm:app-part07-selected-twin-placement}
Inside the selected passive/outgoing branch, the support loading ratio is fixed to \(\rho_\alpha=4/3\), so the branch lies on the symmetric lowest-twin support slice.  The realized coordinate on that slice is \(\varrho_{\rm phys}=2(1-\epsilon)/3\).  The orbit-lock packet is determined by \((\chi_0,\delta_U,Z_W,\epsilon_W,\epsilon_\eta,\Lambda)\) and is exactly blind to the support ratio \(\zeta\).  Therefore support compensation and orbit lock are separate tests that must both be passed by the same completed moving-throat branch.
\end{theorem}

\claimstatus{Stages 223--225 are \StatusExactClosure{} for the support-ratio extraction, selected curve, primitive ranking thresholds, placement formulas, and support/orbit split.  Actual placement of the nonlinear branch and the final value of \(N_Q-1\) remain \StatusOpen{} unless supplied by a completed PDE calculation.}

\section{What Part VII does and does not prove}
\label{sec:app-part07-what-it-does-not-prove}

Part VII proves that the same-charge mixed sector has a much sharper reduced structure than the older placeholder language suggested.  It proves, within the stated closures, that the one-port static mixed bundle cannot be cited as a source of a new long-range attraction; that the linear dynamic mixed bundle cannot be cited as a new radial kernel class; that resonant help is constrained by an exact dispersive/absorptive tradeoff; that the actual weak-axisymmetric branch test is the transported scalar packet \((\Delta_{\rm norm},\Xi_1)\); that the rigid-mouth actual branch diagonalizes in the physical chart \((U,V)\); and that selected support lies on the symmetric lowest-twin slice.

It does not prove that the completed nonlinear moving-throat PDE realizes the favorable point.  In particular, it does not prove any of the following:
\begin{enumerate}[leftmargin=2.2em]
  \item that the true branch has a new same-charge long-range attractive force;
  \item that the true branch supplies a high-quality pole satisfying the dynamic survival window;
  \item that the true branch lands within the transported static \(|\varepsilon\Xi_1|\) ceiling;
  \item that the rigid-mouth assumptions are realized by the full PDE;
  \item that \(U=V=0\) holds on the actual branch;
  \item that selected twin-support placement and outgoing normalization \(N_Q=1\) both hold without additional branch data.
\end{enumerate}

The safe downstream citation is therefore:
\begin{quote}
Inside the declared one-port, weak-axisymmetric, rigid-mouth, and selected-support closures, the same-charge mixed sector reduces to short-range static product kernels, a resonance/linewidth dynamic window, a transported prefactor packet \((\Delta_{\rm norm},\Xi_1)\), the rigid-mouth Cartesian orbit-lock condition \(U=V=0\), and the selected lowest-twin support slice \(\rho_\alpha=4/3\).  Actual branch realization remains a separate PDE calculation.
\end{quote}

For later use, the practical audit recipe is:
\begin{enumerate}[leftmargin=2.2em]
  \item compute the actual one-port bundle data \((D_0,N_0,P_0)\) and check the static product-kernel assumptions;
  \item evaluate the branch-compatible target surface \cref{eq:app-part07-compatibility-equation};
  \item form \((\Delta_{\rm norm},\Xi_1)\) and test the transported ceiling \cref{eq:app-part07-thm-ceiling};
  \item if the branch is rigid-mouth, compute \((U,V)\) from \cref{eq:app-part07-uv-def};
  \item test orbit lock by \(U=V=0\);
  \item compute \(\epsilon\), \(\varrho_{\rm phys}\), the selected ranking region, the support packet, and \(N_Q-1\) using \cref{eq:app-part07-main-final-packet,eq:app-part07-stage225-support-packet}.
\end{enumerate}
This is exactly the point at which Part VIII can either preserve the strict front-end packet or declare a relaxed/open-system companion with an explicit recovery map back to the Stage-225 front end.


\section{Original-stage verification cards}
\label{sec:app-part07-original-stage-verification-cards}

The cards below restore the original Stage~202--225 order.  They are intentionally more compact than the theorem narrative above: each card records the stage purpose, inputs, derivation ledger, boxed output, checks, and downstream use.  The cards are not independent proofs; they are a verification index for the corresponding source stages.



\paragraph{Source layout.}
Stages~202--225 are now sourced from the canonical files \StageFile{stages/stage_202.tex} through \StageFile{stages/stage_225.tex}.  The raw Markdown notes and audit artifacts are tracked in a repo-local provenance index, so they do not have to appear in the reader-facing PDF.  The stage files themselves carry the executable SymPy and Mathematica audit references.

\begingroup
\let\clearpage\relax
\input{stages/stage_202}
\clearpage
\section[Stage 203]{Stage 203: Scalar graph-slice theorem}
\label{stage:203}

\stagefield{Part / anchor}{Part VI; primary anchor \resultanchor{MTDC-T10}, local subanchor \resultanchor{MTDC-T10.1}.}
\claimstatus{\StatusExactClosure{}.}

\stagefield{Purpose}{Stage~203 proves that once a branch is aligned with the target graph, the only remaining closure equation is the scalar condition \(\widehat\chi_Q(\mathbf y)=1\).}

\stagefield{Inputs}{Imports the Stage~202 graph, the graph-error packet, and differentiable candidate families in the positive free-quintuple domain.}

\stagefield{Verification}{SymPy audit: \StageFile{scripts/moving_throat_pde_stage203_free_quintuple_scalar_closure_slice_and_crossing_theorem_sympy_audit.py}.  Mathematica audit: \StageFile{mathematica/moving_throat_pde_stage203_free_quintuple_scalar_closure_slice_and_crossing_theorem_mathematica_audit.wl}.}

\stagefield{Derivation ledger}{The derivation composes \(\chi_Q\) with the target graph, proves graph-lifted tangents lie in \(\ker M_*\), separates graph-aligned and off-graph family packets, and applies the intermediate value theorem to one-parameter graph crossings.}

\stagefield{Output}{Scalar graph slice \(\mathcal Z_*=\{\mathbf x_*^{\rm graph}(\mathbf y):\widehat\chi_Q(\mathbf y)=1\}\), graph-family tangency, and one-parameter crossing theorem.}

\stagefield{Verification note}{The detailed theorem-block formulas supporting this card are collected in the Part VI appendix narrative above and in the source stage.  The card restores original stage order and records the claim boundary; it should not be read as a second independent proof.}

\stagefield{Checks}{\begin{verificationchecklist}
\item Check target monomial ratios are formed against the declared target values before taking logarithms.
\item Check the same-free-quintuple convention is used; changing the free/dependent split changes the projection chart.
\item Check graph-aligned closure is not confused with actual PDE realization.
\item Carry the \StatusExactClosure{} status tag whenever citing the compiler itself.
\end{verificationchecklist}}

\stagefield{Downstream use}{This stage feeds the Part VI realization compiler and finite mixed-ray search.  Downstream papers should cite \resultanchor{MTDC-T10.1} for this local stage range and \resultanchor{MTDC-T10} for the full Part VI compiler.}


\clearpage
\section[Stage 204]{Stage 204: Resonance/linewidth tradeoff}
\label{stage:204}

\stagefield{Part / anchor}{Part VII; primary anchor \resultanchor{MTDC-T11}, local subanchor \resultanchor{MTDC-T11.1}.}
\claimstatus{\StatusExactClosure{}.}

\stagefield{Purpose}{Stage~204 analyzes the only surviving linear dynamic escape hatch: resonant dispersive enhancement of the already-short-range mixed kernels.}

\stagefield{Inputs}{Imports Stage~203 and expands near a simple conservative pole $F_0(\omega_*)=0$ with $F_0^\prime(\omega_*)\ne0$; the wall-like specialization also assumes $D_0(\omega_*)=0$ and $\Delta_0(\omega_*)\ne0$.}

\stagefield{Verification}{SymPy audit: \StageFile{scripts/moving_throat_pde_stage204_resonance_linewidth_tradeoff_dispersive_no_free_lunch_theorem_and_linear_survival_window_sympy_audit.py}.  Mathematica audit: \StageFile{mathematica/moving_throat_pde_stage204_resonance_linewidth_tradeoff_dispersive_no_free_lunch_theorem_and_linear_survival_window_mathematica_audit.wl}.}

\stagefield{Derivation ledger}{The derivation reduces the susceptibility to the Breit--Wigner form $A_*/(\delta-i\gamma_*)$, computes the real and imaginary line shapes, rewrites the result in absorbed-power language, and derives the low-loss survival window.}

\stagefield{Output}{Dispersive/no-free-lunch theorem: $|\Re\chi_*|/|\Im\chi_*|=|\delta|/\gamma_*$, so resonant conservative gain is inseparable from absorptive loading, with maximal balanced leverage at equal conservative and absorptive magnitudes.}

\stagefield{Verification note}{The detailed theorem-block formulas supporting this card are collected in the Part VII appendix narrative above and in the source stage.  The card preserves the original stage order and records the claim boundary; it should not be read as a second independent proof of actual nonlinear branch realization.}

\stagefield{Checks}{\begin{verificationchecklist}
\item Check the stage assumptions against the declared one-port, selected-branch, weak-axisymmetric, or rigid-mouth closure before citing the output.
\item Check whether the row is algebraic, numerical placement, or open branch-realization before transferring it to another paper.
\item Check that same-charge statements are not reinterpreted as a new long-range attraction; all static kernel claims inherit the source-product family limits.
\item Check that orbit-lock claims require the rigid-mouth packet variables and do not follow from generic mouth tracking alone.
\end{verificationchecklist}}

\stagefield{Downstream use}{This stage feeds the Part VII same-charge sieve, rigid-mouth normal form, and selected twin-support branch.  Downstream papers should cite \resultanchor{MTDC-T11.1} for this local stage range and \resultanchor{MTDC-T11} for the full Part VII compiler.}


\clearpage
\input{stages/stage_205}
\clearpage
\input{stages/stage_206}
\clearpage
\section[Stage 207]{Stage 207: Branch-packet compiler and actual-branch kill test}
\label{stage:207}

\stagefield{Part / anchor}{Part VII; primary anchor \resultanchor{MTDC-T11}, local subanchor \resultanchor{MTDC-T11.2}.}
\claimstatus{Mixed: \StatusExactClosure{}, \StatusOpen{}.}

\stagefield{Purpose}{Stage~207 converts the primitive-branch compatibility result into the actual weak-axisymmetric branch packet that the moving-throat PDE must return.}

\stagefield{Inputs}{Imports the normalized prefactor packet $(\Delta_{\rm norm},a_{P_0},b_{P_0})$, the weak-axisymmetric signature $(1,1/2,-1)$, the line $b_{P_0}=3a_{P_0}$, and the target scale $T_{\rm quad}=54Gc_s^5/(5a^5c^5)$.}

\stagefield{Verification}{SymPy audit: \StageFile{scripts/moving_throat_pde_stage207_pde_branch_packet_compiler_weak_axisymmetric_ceiling_transport_and_first_actual_branch_kill_test_sympy_audit.py}.  Mathematica audit: none yet.}

\stagefield{Derivation ledger}{The derivation maps grouped lane prefactors into trace/anomaly coordinates, rewrites the anisotropic packet in terms of $\Xi_1=P_1/P_0$, and transports the static normalization margin through the one-scalar weak-axisymmetric ceiling.}

\stagefield{Output}{Actual-branch kill test: the branch survives only if the transported quantity $(\Delta_{\rm norm}+T_{\rm quad})(1+|\varepsilon\Xi_1|)/\hat m_0^2$ stays below the critical prefactor budget.}

\stagefield{Verification note}{The detailed theorem-block formulas supporting this card are collected in the Part VII appendix narrative above and in the source stage.  The card preserves the original stage order and records the claim boundary; it should not be read as a second independent proof of actual nonlinear branch realization.}

\stagefield{Checks}{\begin{verificationchecklist}
\item Check the stage assumptions against the declared one-port, selected-branch, weak-axisymmetric, or rigid-mouth closure before citing the output.
\item Check whether the row is algebraic, numerical placement, or open branch-realization before transferring it to another paper.
\item Check that same-charge statements are not reinterpreted as a new long-range attraction; all static kernel claims inherit the source-product family limits.
\item Check that orbit-lock claims require the rigid-mouth packet variables and do not follow from generic mouth tracking alone.
\end{verificationchecklist}}

\stagefield{Downstream use}{This stage feeds the Part VII same-charge sieve, rigid-mouth normal form, and selected twin-support branch.  Downstream papers should cite \resultanchor{MTDC-T11.2} for this local stage range and \resultanchor{MTDC-T11} for the full Part VII compiler.}



\clearpage
\section[Stage 208]{Stage 208: Pairwise mixed-ray audit}
\label{stage:208}

\stagefield{Part / anchor}{Part VI; primary anchor \resultanchor{MTDC-T10}, local subanchor \resultanchor{MTDC-T10.3}.}
\claimstatus{Mixed: \StatusExactClosure{}, \StatusNumerical{}.}

\stagefield{Purpose}{Stage~208 opens the first genuinely mixed search sector by allowing two primitive coordinates to move together.}

\stagefield{Inputs}{Imports primitive axis data, off-diagonal Hessian entries, and a pairwise cone direction \(\mathbf s=a_i\widehat{\mathbf e}_i+a_j\widehat{\mathbf e}_j\) with \(a_i^2+a_j^2=1\).}

\stagefield{Verification}{SymPy audit: \StageFile{scripts/moving_throat_pde_stage208_pairwise_mixed_rays_and_off_diagonal_hessian_synergy_sympy_audit.py}.  Mathematica audit: none yet.}

\stagefield{Derivation ledger}{The derivation proves the gradient-synergy law, isolates the off-diagonal curvature leverage \(2a_ia_jh_{ij}\), constructs pairwise curvature envelopes, compares the gradient-optimal and equal-mix screen rays, and states the pairwise promotion rule.}

\stagefield{Output}{Pairwise mixed-ray cone, gradient and off-diagonal Hessian synergy laws, canonical two-ray screen, promotion/deferral rule, and minimal packet for the full ratio optimizer.}

\stagefield{Verification note}{The detailed theorem-block formulas supporting this card are collected in the Part VI appendix narrative above and in the source stage.  The card restores original stage order and records the claim boundary; it should not be read as a second independent proof.}

\stagefield{Checks}{\begin{verificationchecklist}
\item Check all pairwise and triple boundary faces are imported from already-certified lower-support ledgers.
\item Check admissibility windows before including stationary candidates.
\item Check interval comparison uses certified lower/upper endpoints, not just central predictor values.
\item Mark candidate evaluation as \StatusNumerical{} when actual local packet values are inserted.
\end{verificationchecklist}}

\stagefield{Downstream use}{This stage feeds the Part VI realization compiler and finite mixed-ray search.  Downstream papers should cite \resultanchor{MTDC-T10.3} for this local stage range and \resultanchor{MTDC-T10} for the full Part VI compiler.}



\clearpage
\input{stages/stage_209}
\clearpage
\section[Stage 210]{Stage 210: Three-coordinate simplex audit}
\label{stage:210}

\stagefield{Part / anchor}{Part VI; primary anchor \resultanchor{MTDC-T10}, local subanchor \resultanchor{MTDC-T10.3}.}
\claimstatus{Mixed: \StatusExactClosure{}, \StatusNumerical{}.}

\stagefield{Purpose}{Stage~210 promotes the search to genuine three-coordinate positive simplices and proves that their boundaries are already the Stage~209 pairwise ledgers.}

\stagefield{Inputs}{Imports primitive triples, the positive spherical simplex \(a_i^2+a_j^2+a_k^2=1\), pairwise edge packets, and local gradient/Hessian data.}

\stagefield{Verification}{SymPy audit: \StageFile{scripts/moving_throat_pde_stage210_three_coordinate_mixed_simplex_audit_sympy_audit.py}.  Mathematica audit: none yet.}

\stagefield{Derivation ledger}{The derivation reduces every boundary edge to an imported pairwise search, proves the three-coordinate gradient-synergy law, computes total cross leverage, defines fixed-simplex brackets, and gives canonical gradient/equal-mix triple screens.}

\stagefield{Output}{Boundary reduction for triples, three-coordinate gradient and curvature synergy, canonical triple-screen audit, interior-screen dominance criterion, and non-improvement filter.}

\stagefield{Verification note}{The detailed theorem-block formulas supporting this card are collected in the Part VI appendix narrative above and in the source stage.  The card restores original stage order and records the claim boundary; it should not be read as a second independent proof.}

\stagefield{Checks}{\begin{verificationchecklist}
\item Check all pairwise and triple boundary faces are imported from already-certified lower-support ledgers.
\item Check admissibility windows before including stationary candidates.
\item Check interval comparison uses certified lower/upper endpoints, not just central predictor values.
\item Mark candidate evaluation as \StatusNumerical{} when actual local packet values are inserted.
\end{verificationchecklist}}

\stagefield{Downstream use}{This stage feeds the Part VI realization compiler and finite mixed-ray search.  Downstream papers should cite \resultanchor{MTDC-T10.3} for this local stage range and \resultanchor{MTDC-T10} for the full Part VI compiler.}



\clearpage
\section[Stage 211]{Stage 211: Numerator/denominator split and dynamic-window test}
\label{stage:211}

\stagefield{Part / anchor}{Part VII; primary anchor \resultanchor{MTDC-T11}, local subanchor \resultanchor{MTDC-T11.3}.}
\claimstatus{Mixed: \StatusExactClosure{}, \StatusNumerical{}.}

\stagefield{Purpose}{Stage~211 tests the pure-transfer corridor against the dynamic-window machinery.}

\stagefield{Inputs}{Imports the split $\Xi_1=2(\pi_1-\delta_1)$, numerator-rigid and denominator-rigid subcorridors, and the wall-like branch data from the finite-throat primitive slice.}

\stagefield{Verification}{SymPy audit: \StageFile{scripts/moving_throat_pde_stage211_numerator_denominator_split_of_the_pure_transfer_corridor_and_first_actual_dynamic_window_test_sympy_audit.py}.  Mathematica audit: none yet.}

\stagefield{Derivation ledger}{The derivation separates numerator and denominator drift carriers, identifies which subcorridors survive the static constraint, and evaluates the first dynamic-window inequality for the wall-like branch.}

\stagefield{Output}{Numerator/denominator split, dynamic-window audit, and first branch-level comparison between the static $\Xi_1$ budget and resonant dynamic survival.}

\stagefield{Verification note}{The detailed theorem-block formulas supporting this card are collected in the Part VII appendix narrative above and in the source stage.  The card preserves the original stage order and records the claim boundary; it should not be read as a second independent proof of actual nonlinear branch realization.}

\stagefield{Checks}{\begin{verificationchecklist}
\item Check the stage assumptions against the declared one-port, selected-branch, weak-axisymmetric, or rigid-mouth closure before citing the output.
\item Check whether the row is algebraic, numerical placement, or open branch-realization before transferring it to another paper.
\item Check that same-charge statements are not reinterpreted as a new long-range attraction; all static kernel claims inherit the source-product family limits.
\item Check that orbit-lock claims require the rigid-mouth packet variables and do not follow from generic mouth tracking alone.
\end{verificationchecklist}}

\stagefield{Downstream use}{This stage feeds the Part VII same-charge sieve, rigid-mouth normal form, and selected twin-support branch.  Downstream papers should cite \resultanchor{MTDC-T11.3} for this local stage range and \resultanchor{MTDC-T11} for the full Part VII compiler.}



\clearpage
\section[Stage 212]{Stage 212: Selected-branch numerator/denominator signature}
\label{stage:212}

\stagefield{Part / anchor}{Part VII; primary anchor \resultanchor{MTDC-T11}, local subanchor \resultanchor{MTDC-T11.4}.}
\claimstatus{\StatusExactClosure{}.}

\stagefield{Purpose}{Stage~212 pulls the numerator/denominator split onto the selected branch and classifies which side dominates near softening.}

\stagefield{Inputs}{Imports the selected softening-depth variable, the selected branch functions, and the numerator/denominator share map.}

\stagefield{Verification}{SymPy audit: \StageFile{scripts/moving_throat_pde_stage212_selected_branch_numerator_denominator_signature_and_softening_depth_crossover_theorem_sympy_audit.py}.  Mathematica audit: none yet.}

\stagefield{Derivation ledger}{The derivation factorizes the selected branch into numerator-like and denominator-like pieces, derives the classifier threshold $\delta=8/9$, and proves the near-softening denominator-like signature.}

\stagefield{Output}{Selected-branch classifier theorem: the pure-transfer defect is denominator-like near softening, with crossover controlled by the exact threshold $\delta=8/9$.}

\stagefield{Verification note}{The detailed theorem-block formulas supporting this card are collected in the Part VII appendix narrative above and in the source stage.  The card preserves the original stage order and records the claim boundary; it should not be read as a second independent proof of actual nonlinear branch realization.}

\stagefield{Checks}{\begin{verificationchecklist}
\item Check the stage assumptions against the declared one-port, selected-branch, weak-axisymmetric, or rigid-mouth closure before citing the output.
\item Check whether the row is algebraic, numerical placement, or open branch-realization before transferring it to another paper.
\item Check that same-charge statements are not reinterpreted as a new long-range attraction; all static kernel claims inherit the source-product family limits.
\item Check that orbit-lock claims require the rigid-mouth packet variables and do not follow from generic mouth tracking alone.
\end{verificationchecklist}}

\stagefield{Downstream use}{This stage feeds the Part VII same-charge sieve, rigid-mouth normal form, and selected twin-support branch.  Downstream papers should cite \resultanchor{MTDC-T11.4} for this local stage range and \resultanchor{MTDC-T11} for the full Part VII compiler.}



\clearpage
\section[Stage 213]{Stage 213: Dynamic-window classifier and static-first theorem}
\label{stage:213}

\stagefield{Part / anchor}{Part VII; primary anchor \resultanchor{MTDC-T11}, local subanchor \resultanchor{MTDC-T11.4}.}
\claimstatus{Mixed: \StatusExactClosure{}, \StatusNumerical{}.}

\stagefield{Purpose}{Stage~213 translates the selected-branch classifier into a dynamic-window decision rule.}

\stagefield{Inputs}{Imports Stage~212 classifier data, the rigid-split shares, and the survival-window inequality from the resonance theorem.}

\stagefield{Verification}{SymPy audit: \StageFile{scripts/moving_throat_pde_stage213_selected_branch_classifier_to_dynamic_window_compiler_and_static_first_theorem_sympy_audit.py}.  Mathematica audit: none yet.}

\stagefield{Derivation ledger}{The derivation computes the dynamic sign of the selected numerator/denominator split, relates the classifier region to linewidth demand, and compares this with the static prefactor ceiling.}

\stagefield{Output}{Static-first theorem: the static $\Xi_1$ budget is the first kill condition; the dynamic corridor can be checked only after the static transported ceiling is satisfied.}

\stagefield{Verification note}{The detailed theorem-block formulas supporting this card are collected in the Part VII appendix narrative above and in the source stage.  The card preserves the original stage order and records the claim boundary; it should not be read as a second independent proof of actual nonlinear branch realization.}

\stagefield{Checks}{\begin{verificationchecklist}
\item Check the stage assumptions against the declared one-port, selected-branch, weak-axisymmetric, or rigid-mouth closure before citing the output.
\item Check whether the row is algebraic, numerical placement, or open branch-realization before transferring it to another paper.
\item Check that same-charge statements are not reinterpreted as a new long-range attraction; all static kernel claims inherit the source-product family limits.
\item Check that orbit-lock claims require the rigid-mouth packet variables and do not follow from generic mouth tracking alone.
\end{verificationchecklist}}

\stagefield{Downstream use}{This stage feeds the Part VII same-charge sieve, rigid-mouth normal form, and selected twin-support branch.  Downstream papers should cite \resultanchor{MTDC-T11.4} for this local stage range and \resultanchor{MTDC-T11} for the full Part VII compiler.}



\clearpage
\section[Stage 214]{Stage 214: Four-coordinate interior optimizer}
\label{stage:214}

\stagefield{Part / anchor}{Part VI; primary anchor \resultanchor{MTDC-T10}, local subanchor \resultanchor{MTDC-T10.4}.}
\claimstatus{Mixed: \StatusExactClosure{}, \StatusNumerical{}.}

\stagefield{Purpose}{Stage~214 solves the three-parameter interior optimization problem for a primitive quadruple.}

\stagefield{Inputs}{Imports a four-coordinate interior ratio patch, certified root objectives, and compact admissible windows.}

\stagefield{Verification}{SymPy audit: \StageFile{scripts/moving_throat_pde_stage214_full_interior_four_coordinate_simplex_optimizer_and_finite_candidate_reduction_sympy_audit.py}.  Mathematica audit: none yet.}

\stagefield{Derivation ledger}{The derivation forms three stationary numerator equations, introduces one auxiliary square-root variable, obtains a lifted degree pattern \((3,3,3,2)\), proves the \(54\)-candidate preferred bound, and gives a square-root-free projected fallback.}

\stagefield{Output}{Finite quadruple-interior candidate set with preferred bound \(54\) per envelope, projected fallback bound, and interior winner/non-improvement theorem against the Stage~213 boundary ledger.}

\stagefield{Verification note}{The detailed theorem-block formulas supporting this card are collected in the Part VI appendix narrative above and in the source stage.  The card restores original stage order and records the claim boundary; it should not be read as a second independent proof.}

\stagefield{Checks}{\begin{verificationchecklist}
\item Check every proper face of a four- or five-simplex is identified with a previously closed lower-support search.
\item Check lifted algebraic candidates are filtered by positivity, admissibility, and validity intervals.
\item Check support-cardinality closure is limited to the declared five free coordinates.
\item Do not cite the finite sieve as proof that the nonlinear PDE realizes the winning branch.
\end{verificationchecklist}}

\stagefield{Downstream use}{This stage feeds the Part VI realization compiler and finite mixed-ray search.  Downstream papers should cite \resultanchor{MTDC-T10.4} for this local stage range and \resultanchor{MTDC-T10} for the full Part VI compiler.}



\clearpage
\section[Stage 215]{Stage 215: Known 5PN data injection}
\label{stage:215}

\stagefield{Part / anchor}{Part VII; primary anchor \resultanchor{MTDC-T11}, local subanchor \resultanchor{MTDC-T11.4}.}
\claimstatus{Mixed: \StatusNumerical{}, \StatusOpen{}.}

\stagefield{Purpose}{Stage~215 injects the current known Family-1 and 5PN support/source data into the selected-branch inequalities.}

\stagefield{Inputs}{Imports the numerically located support/source packets, the exact selected classifier, and the static-first ceiling.}

\stagefield{Verification}{SymPy audit: \StageFile{scripts/moving_throat_pde_stage215_known_5pn_data_injection_and_current_branch_verdict_sympy_audit.py}.  Mathematica audit: none yet.}

\stagefield{Derivation ledger}{The derivation evaluates the support/source side of the inequalities, compares it with the selected threshold windows, and identifies which gate remains unresolved.}

\stagefield{Output}{Current branch verdict: the support/source side is safe by a large margin in the injected data; the active unresolved gate is static orbit-lock/coherent placement rather than support starvation.}

\stagefield{Verification note}{The detailed theorem-block formulas supporting this card are collected in the Part VII appendix narrative above and in the source stage.  The card preserves the original stage order and records the claim boundary; it should not be read as a second independent proof of actual nonlinear branch realization.}

\stagefield{Checks}{\begin{verificationchecklist}
\item Check the stage assumptions against the declared one-port, selected-branch, weak-axisymmetric, or rigid-mouth closure before citing the output.
\item Check whether the row is algebraic, numerical placement, or open branch-realization before transferring it to another paper.
\item Check that same-charge statements are not reinterpreted as a new long-range attraction; all static kernel claims inherit the source-product family limits.
\item Check that orbit-lock claims require the rigid-mouth packet variables and do not follow from generic mouth tracking alone.
\end{verificationchecklist}}

\stagefield{Downstream use}{This stage feeds the Part VII same-charge sieve, rigid-mouth normal form, and selected twin-support branch.  Downstream papers should cite \resultanchor{MTDC-T11.4} for this local stage range and \resultanchor{MTDC-T11} for the full Part VII compiler.}



\clearpage
\section[Stage 216]{Stage 216: Unique five-coordinate simplex gate}
\label{stage:216}

\stagefield{Part / anchor}{Part VI; primary anchor \resultanchor{MTDC-T10}, local subanchor \resultanchor{MTDC-T10.4}.}
\claimstatus{Mixed: \StatusExactClosure{}, \StatusNumerical{}.}

\stagefield{Purpose}{Stage~216 observes that the five-dimensional free log space has exactly one positive support-five simplex and that all of its faces are support-\(\le4\) searches.}

\stagefield{Inputs}{Imports the support-\(\le4\) ledger, the five primitive axes \(\{\lambda,c,\gamma,U,W\}\), and the unique positive five-simplex.}

\stagefield{Verification}{SymPy audit: \StageFile{scripts/moving_throat_pde_stage216_unique_five_coordinate_positive_simplex_and_support_cardinality_5_gate_sympy_audit.py}.  Mathematica audit: none yet.}

\stagefield{Derivation ledger}{The derivation identifies the boundary faces with Stage~215 quadruple closed simplices, derives the five-coordinate gradient-synergy and ten-way cross-leverage laws, states the five-way canonical screens, and proves the support-cardinality ceiling.}

\stagefield{Output}{Unique five-coordinate simplex, exact face reduction to support-\(\le4\), support-cardinality-five gate, and proof that no higher-support local search exists.}

\stagefield{Verification note}{The detailed theorem-block formulas supporting this card are collected in the Part VI appendix narrative above and in the source stage.  The card restores original stage order and records the claim boundary; it should not be read as a second independent proof.}

\stagefield{Checks}{\begin{verificationchecklist}
\item Check every proper face of a four- or five-simplex is identified with a previously closed lower-support search.
\item Check lifted algebraic candidates are filtered by positivity, admissibility, and validity intervals.
\item Check support-cardinality closure is limited to the declared five free coordinates.
\item Do not cite the finite sieve as proof that the nonlinear PDE realizes the winning branch.
\end{verificationchecklist}}

\stagefield{Downstream use}{This stage feeds the Part VI realization compiler and finite mixed-ray search.  Downstream papers should cite \resultanchor{MTDC-T10.4} for this local stage range and \resultanchor{MTDC-T10} for the full Part VI compiler.}



\clearpage
\section[Stage 217]{Stage 217: Five-coordinate interior optimizer}
\label{stage:217}

\stagefield{Part / anchor}{Part VI; primary anchor \resultanchor{MTDC-T10}, local subanchor \resultanchor{MTDC-T10.4}.}
\claimstatus{Mixed: \StatusExactClosure{}, \StatusNumerical{}.}

\stagefield{Purpose}{Stage~217 solves the last genuine continuum optimization problem: the interior of the unique support-five simplex.}

\stagefield{Inputs}{Imports the support-five ratio chart, certified root objectives, compact admissible windows, and the Stage~216 boundary ledger.}

\stagefield{Verification}{SymPy audit: \StageFile{scripts/moving_throat_pde_stage217_full_interior_five_coordinate_simplex_optimizer_and_finite_candidate_reduction_sympy_audit.py}.  Mathematica audit: none yet.}

\stagefield{Derivation ledger}{The derivation writes the four stationary numerator equations, lifts by one auxiliary square-root variable, obtains degree pattern \((3,3,3,3,2)\), proves the preferred \(179\)-candidate bound per envelope, and gives a square-root-free fallback bound.}

\stagefield{Output}{Finite support-five interior candidate set, preferred budget \(324\), fallback budget \(1500\), and support-five interior improvement/non-improvement theorem.}

\stagefield{Verification note}{The detailed theorem-block formulas supporting this card are collected in the Part VI appendix narrative above and in the source stage.  The card restores original stage order and records the claim boundary; it should not be read as a second independent proof.}

\stagefield{Checks}{\begin{verificationchecklist}
\item Check every proper face of a four- or five-simplex is identified with a previously closed lower-support search.
\item Check lifted algebraic candidates are filtered by positivity, admissibility, and validity intervals.
\item Check support-cardinality closure is limited to the declared five free coordinates.
\item Do not cite the finite sieve as proof that the nonlinear PDE realizes the winning branch.
\end{verificationchecklist}}

\stagefield{Downstream use}{This stage feeds the Part VI realization compiler and finite mixed-ray search.  Downstream papers should cite \resultanchor{MTDC-T10.4} for this local stage range and \resultanchor{MTDC-T10} for the full Part VI compiler.}



\clearpage
\section[Stage 218]{Stage 218: Final local mixed-ray closure theorem}
\label{stage:218}

\stagefield{Part / anchor}{Part VI; primary anchor \resultanchor{MTDC-T10}, local subanchor \resultanchor{MTDC-T10.4}.}
\claimstatus{Mixed: \StatusExactClosure{}, \StatusNumerical{}, \StatusOpen{}.}

\stagefield{Purpose}{Stage~218 splices the unique support-five interior packet to the support-\(\le4\) boundary ledger and closes the declared local mixed-ray search.}

\stagefield{Inputs}{Imports the support-\(\le4\) global ledger, the support-five interior packet, and the unique five-simplex boundary stratification.}

\stagefield{Verification}{SymPy audit: \StageFile{scripts/moving_throat_pde_stage218_full_support_cardinality_5_completion_and_local_mixed_ray_search_closure_sympy_audit.py}.  Mathematica audit: \StageFile{mathematica/moving_throat_pde_stage218_full_support_cardinality_5_completion_and_local_mixed_ray_search_closure_mathematica_audit.wl}.}

\stagefield{Derivation ledger}{The derivation proves the boundary-identification theorem \(\partial\Delta_5^+=\mathcal S_{\le4}^{\rm loc}\), proves the support ceiling, splices the support-five interior and support-\(\le4\) intervals, classifies support-five outcomes, and records the exact evaluation budget.}

\stagefield{Output}{Final local mixed-ray closure theorem, certified support-\(\le5\) splice interval, preferred total budget \(1464\), fallback budget \(2640\), and finite-candidate status of any remaining ambiguity.}

\stagefield{Verification note}{The detailed theorem-block formulas supporting this card are collected in the Part VI appendix narrative above and in the source stage.  The card restores original stage order and records the claim boundary; it should not be read as a second independent proof.}

\stagefield{Checks}{\begin{verificationchecklist}
\item Check every proper face of a four- or five-simplex is identified with a previously closed lower-support search.
\item Check lifted algebraic candidates are filtered by positivity, admissibility, and validity intervals.
\item Check support-cardinality closure is limited to the declared five free coordinates.
\item Do not cite the finite sieve as proof that the nonlinear PDE realizes the winning branch.
\end{verificationchecklist}}

\stagefield{Downstream use}{This stage feeds the Part VI realization compiler and finite mixed-ray search.  Downstream papers should cite \resultanchor{MTDC-T10.4} for this local stage range and \resultanchor{MTDC-T10} for the full Part VI compiler.}

\clearpage
\section[Stage 219]{Stage 219: Microscopic dependent-plane projectors}
\label{stage:219}

\stagefield{Part / anchor}{Part VII; primary anchor \resultanchor{MTDC-T11}, local subanchor \resultanchor{MTDC-T11.6}.}
\claimstatus{\StatusExactClosure{}.}

\stagefield{Purpose}{Stage~219 transfers the direct packet projectors to the microscopic dependent plane.}

\stagefield{Inputs}{Imports the dependent triple $(T_U,K_\eta^{(\rm eff)},\mu_W)$ and the same-free-quintuple graph section.}

\stagefield{Verification}{SymPy audit: \StageFile{scripts/moving_throat_pde_stage219_rigid_mouth_microscopic_dependent_plane_projectors_equal_drift_dressing_ray_and_static_only_restoration_gap_sympy_audit.py}.  Mathematica audit: none yet.}

\stagefield{Derivation ledger}{The derivation computes the dependent-plane compiler, projectors $P_{\rm nt}^{\rm dep}$ and $P_\eta^{\rm dep}$, the equal-drift dressing ray, and the static-only restoration direction.}

\stagefield{Output}{Microscopic projector theorem: clearing the static nontracking defect moves along a different dependent direction than clearing the post-static dressing defect; a static-only fix leaves a dressing gap.}

\stagefield{Verification note}{The detailed theorem-block formulas supporting this card are collected in the Part VII appendix narrative above and in the source stage.  The card preserves the original stage order and records the claim boundary; it should not be read as a second independent proof of actual nonlinear branch realization.}

\stagefield{Checks}{\begin{verificationchecklist}
\item Check the stage assumptions against the declared one-port, selected-branch, weak-axisymmetric, or rigid-mouth closure before citing the output.
\item Check whether the row is algebraic, numerical placement, or open branch-realization before transferring it to another paper.
\item Check that same-charge statements are not reinterpreted as a new long-range attraction; all static kernel claims inherit the source-product family limits.
\item Check that orbit-lock claims require the rigid-mouth packet variables and do not follow from generic mouth tracking alone.
\end{verificationchecklist}}

\stagefield{Downstream use}{This stage feeds the Part VII same-charge sieve, rigid-mouth normal form, and selected twin-support branch.  Downstream papers should cite \resultanchor{MTDC-T11.6} for this local stage range and \resultanchor{MTDC-T11} for the full Part VII compiler.}



\clearpage
\section[Stage 220]{Stage 220: Dynamic mixed-port kernel and phase-lag no-go}
\label{stage:220}

\stagefield{Part / anchor}{Part VII; primary anchor \resultanchor{MTDC-T11}, local subanchor \resultanchor{MTDC-T11.1}.}
\claimstatus{\StatusExactClosure{}.}

\stagefield{Purpose}{Stage~220 promotes the one-port bundle to finite frequency and tests whether linear dynamic driving changes the same-charge kernel class.}

\stagefield{Inputs}{Imports Stage~219, frequency-dependent BdG and Maxwell/mixed denominators, the outgoing dressing $\Pi_{\rm out}(\omega)$, and the off-pole conditions $\varpi^2-\omega^2\ne0$, $\Delta_\Pi\ne0$, $D_\Pi\ne0$.}

\stagefield{Verification}{SymPy audit: \StageFile{scripts/moving_throat_pde_stage220_dynamic_mixed_port_kernel_phase_lag_no_go_and_resonant_survival_gate_sympy_audit.py}.  Mathematica audit: none yet.}

\stagefield{Derivation ledger}{The derivation builds the dynamic $3\times3$ susceptibility, proves the determinant identity $\det\mathcal K_\Pi=\Delta_\Pi D_\Pi$, repeats the product-family factorization at finite frequency, differentiates the outgoing-port insertion, and separates in-phase conservative reshaping from quadrature pumping.}

\stagefield{Output}{Dynamic product-family theorem and phase-lag no-go: linear dynamic mixing preserves the static spatial families, and the first-order passive/outgoing insertion is a phase-lag term rather than a new real barrier-lowering kernel off resonance.}

\stagefield{Verification note}{The detailed theorem-block formulas supporting this card are collected in the Part VII appendix narrative above and in the source stage.  The card preserves the original stage order and records the claim boundary; it should not be read as a second independent proof of actual nonlinear branch realization.}

\stagefield{Checks}{\begin{verificationchecklist}
\item Check the stage assumptions against the declared one-port, selected-branch, weak-axisymmetric, or rigid-mouth closure before citing the output.
\item Check whether the row is algebraic, numerical placement, or open branch-realization before transferring it to another paper.
\item Check that same-charge statements are not reinterpreted as a new long-range attraction; all static kernel claims inherit the source-product family limits.
\item Check that orbit-lock claims require the rigid-mouth packet variables and do not follow from generic mouth tracking alone.
\end{verificationchecklist}}

\stagefield{Downstream use}{This stage feeds the Part VII same-charge sieve, rigid-mouth normal form, and selected twin-support branch.  Downstream papers should cite \resultanchor{MTDC-T11.1} for this local stage range and \resultanchor{MTDC-T11} for the full Part VII compiler.}



\clearpage
\section[Stage 221]{Stage 221: Physical transfer-shape compiler}
\label{stage:221}

\stagefield{Part / anchor}{Part VII; primary anchor \resultanchor{MTDC-T11}, local subanchor \resultanchor{MTDC-T11.6}.}
\claimstatus{\StatusExactClosure{}.}

\stagefield{Purpose}{Stage~221 rewrites the rigid-mouth packet in the physical transfer-shape variables.}

\stagefield{Inputs}{Imports $\mathcal T^2$, $R_{\rm target}$, $\epsilon_\eta$, and the constant $\Lambda_0$.}

\stagefield{Verification}{SymPy audit: \StageFile{scripts/moving_throat_pde_stage221_physical_branch_transfer_shape_compiler_packet_factorization_and_post_static_dressing_invariance_theorem_sympy_audit.py}.  Mathematica audit: none yet.}

\stagefield{Derivation ledger}{The derivation proves $R_{\rm target}\mathcal T^2=\Lambda_0(1-\epsilon_\eta)$, factorizes the finite packet, and splits the branch gates into tracking, transfer-shape, and dressing conditions.}

\stagefield{Output}{Physical transfer-shape compiler: the finite rigid-mouth packet is completely described by tracking, $\mathcal T^2$, and $\epsilon_\eta$.}

\stagefield{Verification note}{The detailed theorem-block formulas supporting this card are collected in the Part VII appendix narrative above and in the source stage.  The card preserves the original stage order and records the claim boundary; it should not be read as a second independent proof of actual nonlinear branch realization.}

\stagefield{Checks}{\begin{verificationchecklist}
\item Check the stage assumptions against the declared one-port, selected-branch, weak-axisymmetric, or rigid-mouth closure before citing the output.
\item Check whether the row is algebraic, numerical placement, or open branch-realization before transferring it to another paper.
\item Check that same-charge statements are not reinterpreted as a new long-range attraction; all static kernel claims inherit the source-product family limits.
\item Check that orbit-lock claims require the rigid-mouth packet variables and do not follow from generic mouth tracking alone.
\end{verificationchecklist}}

\stagefield{Downstream use}{This stage feeds the Part VII same-charge sieve, rigid-mouth normal form, and selected twin-support branch.  Downstream papers should cite \resultanchor{MTDC-T11.6} for this local stage range and \resultanchor{MTDC-T11} for the full Part VII compiler.}



\clearpage
\section[Stage 222]{Stage 222: Concrete finite-throat primitive branch}
\label{stage:222}

\stagefield{Part / anchor}{Part VII; primary anchor \resultanchor{MTDC-T11}, local subanchor \resultanchor{MTDC-T11.2}.}
\claimstatus{Mixed: \StatusExactClosure{}, \StatusNumerical{}.}

\stagefield{Purpose}{Stage~222 inserts the lowest concrete finite-throat primitive branch into the resonance/linewidth theorem.}

\stagefield{Inputs}{Imports the lowest N/N wall--$U$ profile, the lowest D/N half-wave support--$W$ profile, the one-port mixed kernel, and the line-shape survival gate from Stage~221.}

\stagefield{Verification}{SymPy audit: \StageFile{scripts/moving_throat_pde_stage222_concrete_finite_throat_primitive_branch_pole_census_and_residue_linewidth_survival_test_sympy_audit.py}.  Mathematica audit: none yet.}

\stagefield{Derivation ledger}{The derivation evaluates the exact overlap constants for the primitive branch, forms the quartic pole polynomial, identifies wall-like and internal poles, computes residues and linewidths, and tests the resulting low-loss survival inequality on an illustrative admissible slice.}

\stagefield{Output}{Concrete branch pole census, residue/linewidth cancellation, and first branch-level dynamic survival test.  The pole polynomial and survival gate are exact in the primitive closure; the displayed slice is a numerical placement.}

\stagefield{Verification note}{The detailed theorem-block formulas supporting this card are collected in the Part VII appendix narrative above and in the source stage.  The card preserves the original stage order and records the claim boundary; it should not be read as a second independent proof of actual nonlinear branch realization.}

\stagefield{Checks}{\begin{verificationchecklist}
\item Check the stage assumptions against the declared one-port, selected-branch, weak-axisymmetric, or rigid-mouth closure before citing the output.
\item Check whether the row is algebraic, numerical placement, or open branch-realization before transferring it to another paper.
\item Check that same-charge statements are not reinterpreted as a new long-range attraction; all static kernel claims inherit the source-product family limits.
\item Check that orbit-lock claims require the rigid-mouth packet variables and do not follow from generic mouth tracking alone.
\end{verificationchecklist}}

\stagefield{Downstream use}{This stage feeds the Part VII same-charge sieve, rigid-mouth normal form, and selected twin-support branch.  Downstream papers should cite \resultanchor{MTDC-T11.2} for this local stage range and \resultanchor{MTDC-T11} for the full Part VII compiler.}

\clearpage
\section[Stage 223]{Stage 223: Isotropic target surface and primitive compatibility}
\label{stage:223}

\stagefield{Part / anchor}{Part VII; primary anchor \resultanchor{MTDC-T11}, local subanchor \resultanchor{MTDC-T11.2}.}
\claimstatus{Mixed: \StatusExactClosure{}, \StatusNumerical{}.}

\stagefield{Purpose}{Stage~223 asks whether the same primitive branch can satisfy the conservative one-pole condition and the outgoing normalization target simultaneously.}

\stagefield{Inputs}{Imports the isotropic full-bundle moments $D_0,D_2,D_4,N_0,N_2,N_4$, the one-pole identity $u_4=4u_2^2$, and the universal target $P_{0,{\rm target}}$.}

\stagefield{Verification}{SymPy audit: \StageFile{scripts/moving_throat_pde_stage223_5pn_isotropic_target_surface_primitive_branch_compatibility_and_dynamic_survival_window_sympy_audit.py}.  Mathematica audit: none yet.}

\stagefield{Derivation ledger}{The derivation eliminates the static stiffness through the one-pole condition and compares the resulting $D_0$ with the normalization equation $P_0=N_0/D_0=P_{0,{\rm target}}$.}

\stagefield{Output}{Exact primitive compatibility surface $N_0/P_{0,{\rm target}}=3(M+B_2+Z_2)^2/(B_4+Z_4)$ and finite dynamic survival windows on the compatible branch.}

\stagefield{Verification note}{The detailed theorem-block formulas supporting this card are collected in the Part VII appendix narrative above and in the source stage.  The card preserves the original stage order and records the claim boundary; it should not be read as a second independent proof of actual nonlinear branch realization.}

\stagefield{Checks}{\begin{verificationchecklist}
\item Check the stage assumptions against the declared one-port, selected-branch, weak-axisymmetric, or rigid-mouth closure before citing the output.
\item Check whether the row is algebraic, numerical placement, or open branch-realization before transferring it to another paper.
\item Check that same-charge statements are not reinterpreted as a new long-range attraction; all static kernel claims inherit the source-product family limits.
\item Check that orbit-lock claims require the rigid-mouth packet variables and do not follow from generic mouth tracking alone.
\end{verificationchecklist}}

\stagefield{Downstream use}{This stage feeds the Part VII same-charge sieve, rigid-mouth normal form, and selected twin-support branch.  Downstream papers should cite \resultanchor{MTDC-T11.2} for this local stage range and \resultanchor{MTDC-T11} for the full Part VII compiler.}



\clearpage
\section[Stage 224]{Stage 224: Primitive ranking on selected twin-support branch}
\label{stage:224}

\stagefield{Part / anchor}{Part VII; primary anchor \resultanchor{MTDC-T11}, local subanchor \resultanchor{MTDC-T11.7}.}
\claimstatus{\StatusExactClosure{}.}

\stagefield{Purpose}{Stage~224 ranks the primitive carrier weights along the selected twin-support curve.}

\stagefield{Inputs}{Imports $\epsilon_*=1-3\varrho/2$, $\sigma=4/(3\varrho)-2$, $0<\varrho<2/3$, and the constructive coherent bound $0<\beta<2/11$.}

\stagefield{Verification}{SymPy audit: \StageFile{scripts/moving_throat_pde_stage224_exact_primitive_ranking_on_the_selected_twin_support_branch_sympy_audit.py}.  Mathematica audit: none yet.}

\stagefield{Derivation ledger}{The derivation computes the primitive coherent weights, proves $w_\chi=2w_Z$, derives the thresholds $\varrho_{W\Lambda}$ and $\varrho_{U\Lambda}$, and orders the three ranking regions.}

\stagefield{Output}{Primitive ranking theorem: the selected branch has a complete phase diagram with thresholds $\varrho_{W\Lambda}=2(1+\beta^2)/[3(2+\beta^2)]$ and $\varrho_{U\Lambda}=2(1+\beta^2)/[3(1+\beta+\beta^2)]$.}

\stagefield{Verification note}{The detailed theorem-block formulas supporting this card are collected in the Part VII appendix narrative above and in the source stage.  The card preserves the original stage order and records the claim boundary; it should not be read as a second independent proof of actual nonlinear branch realization.}

\stagefield{Checks}{\begin{verificationchecklist}
\item Check the stage assumptions against the declared one-port, selected-branch, weak-axisymmetric, or rigid-mouth closure before citing the output.
\item Check whether the row is algebraic, numerical placement, or open branch-realization before transferring it to another paper.
\item Check that same-charge statements are not reinterpreted as a new long-range attraction; all static kernel claims inherit the source-product family limits.
\item Check that orbit-lock claims require the rigid-mouth packet variables and do not follow from generic mouth tracking alone.
\end{verificationchecklist}}

\stagefield{Downstream use}{This stage feeds the Part VII same-charge sieve, rigid-mouth normal form, and selected twin-support branch.  Downstream papers should cite \resultanchor{MTDC-T11.7} for this local stage range and \resultanchor{MTDC-T11} for the full Part VII compiler.}



\clearpage
\section[Stage 225]{Stage 225: Actual twin-support placement and coherent orbit-lock compiler}
\label{stage:225}

\stagefield{Part / anchor}{Part VII; primary anchor \resultanchor{MTDC-T11}, local subanchor \resultanchor{MTDC-T11.7}.}
\claimstatus{Mixed: \StatusExactClosure{}, \StatusOpen{}.}

\stagefield{Purpose}{Stage~225 places the actual coherent branch on the selected twin-support curve and packages the final front-end branch-evaluation data.}

\stagefield{Inputs}{Imports the selected curve from Stage~224, the coherent local D/N placement map, the physical variable $\epsilon$, and the support-blind orbit-lock packet compiler.}

\stagefield{Verification}{SymPy audit: \StageFile{scripts/moving_throat_pde_stage225_actual_twin_support_placement_and_coherent_orbit_lock_compiler_sympy_audit.py}.  Mathematica audit: \StageFile{mathematica/moving_throat_pde_stage225_actual_twin_support_placement_and_coherent_orbit_lock_compiler_mathematica_audit.wl}.}

\stagefield{Derivation ledger}{The derivation identifies $\varrho_{\rm phys}=2(1-\epsilon)/3$, rewrites the thresholds in the realized variable, classifies the support lane, constructs the finite orbit packet, and separates support placement from orbit lock.}

\stagefield{Output}{Final front-end packet: actual selected placement coordinate, primitive ranking region, support/orbit packet split, weak-axisymmetric orbit packet, and separate outgoing finish-line datum $N_Q-1$.}

\stagefield{Verification note}{The detailed theorem-block formulas supporting this card are collected in the Part VII appendix narrative above and in the source stage.  The card preserves the original stage order and records the claim boundary; it should not be read as a second independent proof of actual nonlinear branch realization.}

\stagefield{Checks}{\begin{verificationchecklist}
\item Check the stage assumptions against the declared one-port, selected-branch, weak-axisymmetric, or rigid-mouth closure before citing the output.
\item Check whether the row is algebraic, numerical placement, or open branch-realization before transferring it to another paper.
\item Check that same-charge statements are not reinterpreted as a new long-range attraction; all static kernel claims inherit the source-product family limits.
\item Check that orbit-lock claims require the rigid-mouth packet variables and do not follow from generic mouth tracking alone.
\end{verificationchecklist}}

\stagefield{Downstream use}{This stage feeds the Part VII same-charge sieve, rigid-mouth normal form, and selected twin-support branch.  Downstream papers should cite \resultanchor{MTDC-T11.7} for this local stage range and \resultanchor{MTDC-T11} for the full Part VII compiler.}

\endgroup
